\subsection{Transitivity of Reduced Norms and Traces}

PROPOSITION 7. --- Let L be a maximal commutative semisimple subalgebra of A and a an element of L. We have

\[
\mathrm{Pcrd} _ {\mathrm{A/K}} (a; \mathrm{X}) = \mathrm{Pc} _ {\mathrm{L/K}} (a; \mathrm{X}),\tag{52}
\]

\[
\mathrm{Trd} _ {\mathrm{A/K}} (a) = \mathrm{Tr} _ {\mathrm{L/K}} (a),\tag{53}
\]

\[
\mathrm{Nrd} _ {\mathrm{A/K}} (a) = \mathrm{N} _ {\mathrm{L/K}} (a).\tag{54}
\]

By Proposition 3 of VIII, p.~262, the L-modules A and \(L^{n}\) are isomorphic; we therefore have the relation

\[
\mathrm{Pc} _ {\mathrm{A/K}} (a; \mathrm{X}) = \mathrm{Pc} _ {\mathrm{L/K}} (a; \mathrm{X}) ^ {n}.
\]

Since the polynomial \(\mathrm{Pc}_{\mathrm{L}/\mathrm{K}}(a;\mathrm{X})\) is monic, it is therefore equal to the reduced characteristic polynomial \(\mathrm{Pcrd}_{\mathrm{A}/\mathrm{K}}(a;\mathrm{X})\) (VIII, p.~339, Lemma 2); this gives formula (52). By comparing the coefficients of \(X^{n-1}\) (resp. the constant terms) on each side of (52), we obtain formula (53) (resp. (54)).

COROLLARY. --- Let D be a field of finite degree over K with center K. Let a be an element of K and L a maximal commutative subfield of D containing a. We have

\[
\begin{array}{c} \mathrm {Pcrd_ {D / K}} (a; \mathrm{X}) = \mathrm {Pc_ {L / K}} (a; \mathrm{X}), \\ \mathrm{Tr} _ {\mathrm{D/K}} (a) = \mathrm{Tr} _ {\mathrm{L/K}} (a), \\ \mathrm{Nrd} _ {\mathrm{D/K}} (a) = \mathrm{N} _ {\mathrm{L/K}} (a). \end{array}\tag{55}
\]

Indeed, a maximal commutative subfield L of D is a maximal commutative semisimple subalgebra of D by Corollary 2 of VIII, p.~265.

PROPOSITION 8. --- Let B be a simple subalgebra of A. Denote the center of B by L and the commutant of B in A by B'. Then B' is a central simple algebra over the field L; we denote its reduced degree by r. For every element b of B, we have the relations

\[
\mathrm{Pcrd} _ {\mathrm{A/K}} (b; \mathrm{X}) = \mathrm{N} _ {\mathrm{L[X]/K[X]}} (\mathrm{Pcrd} _ {\mathrm{B/L}} (b; \mathrm{X})) ^ {r},\tag{56}
\]

\[
\mathrm{Trd} _ {\mathrm{A/K}} (b) = r \mathrm{Tr} _ {\mathrm{L/K}} (\mathrm{Trd} _ {\mathrm{B/K}} (b)),\tag{57}
\]

\[
\mathrm{Nrd} _ {\mathrm{A/K}} (b) = \mathrm{N} _ {\mathrm{L/K}} (\mathrm{Nrd} _ {\mathrm{B/L}} (b)) ^ {r}.\tag{58}
\]

Lemma 6. --- Let \(\mathbf{K}'\) be a commutative algebra of finite degree \(d\) over \(\mathbf{K}\) and

\[
\mathrm{P} (\mathrm{X}) = \mathrm{X} ^ {s} + a _ {1} \mathrm{X} ^ {s - 1} + \dots + a _ {s}
\]

a monic polynomial with coefficients in \(K'\) . Then the polynomial \(Q = N_{K'[X]/K[X]}(P)\) in \(K[X]\) is monic of degree sd, the coefficient of \(X^{sd-1}\) in \(Q(X)\) is equal to \(\mathrm{Tr}_{\mathrm{K}'/\mathrm{K}}(a_1)\) , and the constant term of Q is \(N_{K'/K}(a_s)\) .

We denote the \(K^{\prime}\) -algebra \(\mathrm{K}^{\prime}[\mathrm{T}]/( \mathrm{P}(\mathrm{T}))\) by \(K^{\prime\prime}\) and the canonical class of T in \(K^{\prime\prime}\) by t. The sequence \((1, t, \ldots, t^{s-1})\) is a basis of \(K^{\prime\prime}\) over \(K^{\prime}\) , and the matrix of multiplication by t in this basis is of the form

\[
\tau = \left( \begin{array}{c c c c c} 0 & 0 & \dots & 0 & - a _ {s} \\ 1 & 0 & \dots & 0 & - a _ {s - 1} \\ 0 & 1 & \dots & 0 & - a _ {s - 2} \\ . & . & \dots & . & . \\ . & . & \dots & . & . \\ 0 & . & \dots & 1 & - a _ {1} \end{array} \right).\tag{59}
\]

The determinant of \(XI_{n}-\tau\) is calculated by induction on s, by expanding along the first row. We obtain \(\det(\mathrm{X}I_{n}-\tau)=\mathrm{P}(\mathrm{X})\) . In other words, we have \(\mathrm{P}(\mathrm{X})=\mathrm{P}\mathrm{c}_{\mathrm{K}^{\prime\prime}/\mathrm{K}^{\prime}}(t;\mathrm{X})\) . In particular, \(\mathrm{Tr}_{\mathrm{K}^{\prime\prime}/\mathrm{K}^{\prime}}(t)=-a_{1}\) and \(\mathrm{N}_{\mathrm{K}^{\prime\prime}/\mathrm{K}^{\prime}}(t)=(-1)^{s}a_{s}\) . By the transitivity formula (III, §9, No.~4, p.~548, Corollary), we have

\[
\begin{array}{r} \mathrm{Tr} _ {\mathrm{K} ^ {\prime \prime} / \mathrm{K}} (t) = - \mathrm{Tr} _ {\mathrm{K} ^ {\prime} / \mathrm{K}} (a _ {1}), \quad \mathrm{N} _ {\mathrm{K} ^ {\prime \prime} / \mathrm{K}} (t) = (- 1) ^ {s d} \mathrm{N} _ {\mathrm{K} ^ {\prime} / \mathrm{K}} (a _ {s}), \\ \mathrm{Q} (\mathrm{X}) = \mathrm{Pc} _ {\mathrm{K} ^ {\prime \prime} / \mathrm{K}} (t; \mathrm{X}). \end{array}
\]

On the other hand, \([K'' : K] = [K'' : K'][K' : K] = sd\) , so \(Q(X)\) is a monic polynomial of degree sd. Lemma 6 follows.

Let us prove Proposition 8. Since the ring B is simple, its center L is a field (VIII, p.~121, Corollary 1). By Theorem 5 of VIII, p.~259, the commutant \(B'\) of B in A is a simple ring with center L, and we have the equality \([A:K]=[B:K][B':K]\) . We denote the reduced degree of \(B'\) over L by r, that of B over L by s, and the degree of L over K by d.~We have

\[
[ \mathrm{A}: \mathrm{K} ] = n ^ {2}, [ \mathrm{B} ^ {\prime}: \mathrm{K} ] = r ^ {2} d, [ \mathrm{B}: \mathrm{K} ] = s ^ {2} d,
\]

and therefore \(n^{2}=r^{2}s^{2}d^{2}\) , that is, n=rsd.

Let \(b\) be an element of B, and let \(\mathrm{P}(\mathrm{X})\) be its reduced characteristic polynomial over the L-algebra B; it is monic of degree \(s\) . By Lemma 6, the polynomial \(\mathrm{Q} = \mathrm{N}_{\mathrm{L}[\mathrm{X}] / \mathrm{K}[\mathrm{X}]}(\mathrm{P})\) is monic of degree \(sd\) . The polynomial \(\mathrm{R} = \mathrm{Q}^r\) is therefore monic of degree \(rsd = n\) . Again by Lemma 6, the coefficient of \(\mathrm{X}^{n - 1}\) in \(\mathrm{R}(\mathrm{X})\) is equal to \(-r\operatorname{Tr}_{\mathrm{L} / \mathrm{K}}(\operatorname{Trd}_{\mathrm{B} / \mathrm{L}}(b))\) , and the constant term of \(\mathrm{R}(\mathrm{X})\) is \((\mathrm{N}_{\mathrm{L} / \mathrm{K}}((-1)^s\mathrm{Nrd}_{\mathrm{B} / \mathrm{L}}(b)))^r = (-1)^n\mathrm{N}_{\mathrm{L} / \mathrm{K}}(\mathrm{Nrd}_{\mathrm{B} / \mathrm{L}}(b))^r\) .

Since \([A:K]=r^{2}d[B:K]\) , the left B-module A is free of rank \(r^{2}d\) (VIII, p.~124, Proposition 5). We therefore have

\[
\mathrm{Pc} _ {\mathrm{A/K}} (b; \mathrm{X}) = \mathrm{Pc} _ {\mathrm{B/K}} (b; \mathrm{X}) ^ {d r ^ {2}}.\tag{60}
\]

By the corollary of III, §9, No.~4, p.~548, we have

\[
\mathrm{Pc} _ {\mathrm{B/K}} (b; \mathrm{X}) = \mathrm{N} _ {\mathrm{L[X] / K[X]}} (\mathrm{Pc} _ {\mathrm{B/L}} (b; \mathrm{X})),\tag{61}
\]

and since \(P(X)\) is the reduced characteristic polynomial of b over the L-algebra B, we have

\[
\mathrm{Pc} _ {\mathrm{B/L}} (b; \mathrm{X}) = \mathrm{P(X)} ^ {s}.\tag{62}
\]

Finally, by formulas (60)---(62) and the definition of R(X), we have

\[
\mathrm{Pc} _ {\mathrm{A/K}} (b; \mathrm{X}) = \mathrm{N} _ {\mathrm{L[X]/K[X]}} (\mathrm{P(X)}) ^ {d r ^ {2} s} = \mathrm{Q(X)} ^ {d r ^ {2} s} = \mathrm{R(X)} ^ {r s d} = \mathrm{R(X)} ^ {n}\tag{63}
\]

so \(\mathrm{R}(\mathrm{X})\) is the reduced characteristic polynomial of b over the K-algebra A.

We have proved formula (56). Formulas (57) and (58) follow immediately from formula (56) and Lemma 6 because the coefficient of \(\mathrm{X}^{n - 1}\) in \(\mathrm{Pcrd}_{\mathrm{A / K}}(b;\mathrm{X})\) is equal to \(-\mathrm{Trd}_{\mathrm{A / K}}(b)\) and the constant term is \((-1)^n\mathrm{Nrd}_{\mathrm{A / K}}(b)\) .

\subsection{Reduced Norms and Determinants}

In this subsection, D is a field of finite degree over K with center K. We denote by \(D_{ab}^{*}\) the quotient of the multiplicative group \(D^{*}\) by its derived (or commutator) group and by \(\pi\) the canonical homomorphism from \(D^{*}\) to \(D_{ab}^{*}\) . The mapping \(Nrd_{D/K}\) induces a group homomorphism from \(D^{*}\) to \(K^{*}\) ; the kernel of this homomorphism contains the derived group of \(D^{*}\) because K is commutative. Hence there exists a unique homomorphism \(\overline{Nrd}\) from \(D_{ab}^{*}\) to \(K^{*}\) such that \(\mathrm{Nrd}_{\mathrm{D/K}}(x)=\overline{\mathrm{Nrd}}\circ\pi(x)\) for every \(x\in D^{*}\) .

PROPOSITION 9. --- Let V be a finite-dimensional right vector space over the field D. Let E be the algebra \(\operatorname{End}_{\mathrm{D}}(\mathrm{V})\) over the field K; it is central simple and of finite degree. For every invertible element u of E, we have

\[
\mathrm{Nrd} _ {\mathrm{E/K}} (u) = \overline {{\mathrm{Nrd}}} (\det u)\tag{64}
\]

(cf.~VIII, p.~452, Proposition 2 for the definition of the determinant \(\det u\) of \(u\) ).

We denote the dimension of V over D by n and identify E with the matrix algebra \(\mathbf{M}_{n}(\mathrm{D})\) using a basis of V over D. The multiplicative group \(\mathrm{GL}_{n}(\mathrm{D})\) of the algebra E is generated by the diagonal matrices and the matrices \(\mathrm{B}_{ij}(\lambda)\) (II, §10, No.~13, p.~362, Corollary 1). Proposition 9 therefore follows from the two specific cases below.

\begin{enumerate}
\def\labelenumi{\Alph{enumi})}
\tightlist
\item
  Suppose that u is the diagonal matrix \(\mathrm{diag}(a_{1},\ldots,a_{n})\) . For every \(1\leqslant i\leqslant n\) , let \(L_{i}\) be a maximal commutative subfield of D containing \(a_{i}\) ; let L be the subalgebra of E consisting of the diagonal matrices \(\mathrm{diag}(t_{1},\ldots,t_{n})\) with \(t_{i}\in L_{i}\) for \(1\leqslant i\leqslant n\) . Let d be the reduced degree of D over K. We have \([L_{i}:K]=d\) for \(1\leqslant i\leqslant n\) (VIII, p.~265, Corollary 2). The K-algebra L is isomorphic to \(L_{1}\times\cdots\times L_{n}\) and therefore semisimple of degree nd; now, we have \([E:K]=n^{2}[D:K]=n^{2}d^{2}=[L:K]^{2}\) . It follows that L is a maximal commutative semisimple subalgebra of E (VIII, p.~262, Proposition 3). By Proposition 7 of VIII, p.~346, we have \(\mathrm{Nrd}_{\mathrm{E}/\mathrm{K}}(u)=\mathrm{N}_{\mathrm{L}/\mathrm{K}}(u)\) . Consequently, by formula (18) of III, §9, No.~3, p.~544, we have
\end{enumerate}

\[
\begin{array}{c} \mathrm{Nrd} _ {\mathrm{E/K}} (u) = \prod_ {i = 1} ^ {n} \mathrm{N} _ {\mathrm{L} _ {i} / \mathrm{K}} (a _ {i}) = \prod_ {i = 1} ^ {n} \mathrm{Nrd} _ {\mathrm{D/K}} (a _ {i}) \\ = \mathrm{Nrd} _ {\mathrm{D/K}} (a _ {1} \dots a _ {n}) = \overline {{\mathrm{Nrd}}} (\pi (a _ {1} \dots a _ {n}))  . \end{array}
\]

Moreover, we have \(\det u = \pi(a_{1} \cdots a_{n})\) by Proposition 3 of VIII, p.~453, which gives formula (64) in this case.

\begin{enumerate}
\def\labelenumi{\Alph{enumi})}
\setcounter{enumi}{1}
\tightlist
\item
  Suppose that u is equal to \(\mathrm{B}_{ij}(\lambda)\) , where \(\lambda\) is an element of D and i,j are distinct integers in the interval \([1,n]\) . Denote by d the reduced degree of D over K and by M the vector space over K deduced from V by restriction of scalars from D to K. Then M is a simple E-module, and we have
\end{enumerate}

\[
\mathrm{Pc} _ {\mathrm{M/K}} (u; \mathrm{X}) = \mathrm{Pcrd} _ {\mathrm{E/K}} (u; \mathrm{X}) ^ {d}\tag{65}
\]

by formula (26) of VIII, p.~341. Moreover, M is a vector space of dimension \(nd^2\) over K, and \(u - 1_{\mathrm{M}}\) is a nilpotent endomorphism of M; we therefore have

\[
\mathrm{Pc} _ {\mathrm{M/K}} (u; \mathrm{X}) = (\mathrm{X} - 1) ^ {n d ^ {2}}.\tag{66}
\]

By comparing formulas (65) and (66), we obtain

\[
\mathrm{Pcrd} _ {\mathrm{E/K}} (u; \mathrm{X}) = (\mathrm{X} - 1) ^ {n d}\tag{67}
\]

and, in particular, \(\mathrm{Nrd}_{\mathrm{E}/\mathrm{K}}(u)=1\) . We also have det u=1 by Proposition 3 of VIII, p.~453; formula (64) therefore holds in this case.

Remark. --- We have \(\mathrm{Nrd}_{\mathrm{E}/\mathrm{K}}(u)=0\) if the element u of E is not invertible (VIII, p.~340, Proposition 3).

\BourbakiExercises{Exercises}

\begin{enumerate}
\def\labelenumi{\arabic{enumi})}
\tightlist
\item
  Let B be an algebra of finite degree over the commutative field K, and let b be an element of B.
\end{enumerate}

\begin{enumerate}
\def\labelenumi{\alph{enumi})}
\item
  If \(b\) is nilpotent, then we have \(\mathrm{Tr}_{\mathrm{M / K}}(b) = 0\) for every B-module M of finite dimension over K. Deduce that if \(b\) belongs to the radical of B, then we have \(\mathrm{Tr}_{\mathrm{M / K}}(bc) = 0\) for every \(c \in \mathrm{B}\) .
\item
  Prove that if \(\operatorname{Tr}_{\mathrm{S}/\mathrm{K}}(bc)=0\) for every simple B-module S and every \(c\in B\) , then b belongs to the radical of B (use the density theorem).
\end{enumerate}

\begin{enumerate}
\def\labelenumi{\arabic{enumi})}
\setcounter{enumi}{1}
\tightlist
\item
  Let B be an algebra of finite degree n over the commutative field K. For any element b of B, denote by \(\mathrm{Pm}_{\mathrm{B}/\mathrm{K}}(b;\mathrm{X})\) the minimal polynomial of b over K (V, §3, No.~1, p.~16). We have \(\mathrm{Pm}_{\mathrm{B}/\mathrm{K}}(b;\mathrm{X})=\mathrm{Pm}_{\mathrm{B}^{\circ}/\mathrm{K}}(b;\mathrm{X})\) .
\end{enumerate}

\begin{enumerate}
\def\labelenumi{\alph{enumi})}
\item
  Prove that \(\mathrm{Pm}_{\mathrm{B / K}}(b;\mathrm{X})\) divides \(\mathrm{Pc}_{\mathrm{B / K}}(b;\mathrm{X})\) , which divides \((\mathrm{Pm}_{\mathrm{B / K}}(b;\mathrm{X}))^n\) .
\item
  For every extension L of K, we have \(\mathrm{Pm}_{\mathrm{B}_{(\mathrm{L})}/\mathrm{L}}(1 \otimes b; \mathrm{X}) = \mathrm{Pm}_{\mathrm{B}/\mathrm{K}}(b; \mathrm{X})\) .
\item
  Let \((e_i)_{i\in \mathrm{I}}\) be a basis of B over K and \(\mathbf{Y} = (\mathrm{Y}_i)_{i\in \mathrm{I}}\) a family of variables. The principal polynomial of the algebra B with respect to the basis \((e_i)\) , denoted by \(\mathrm{P}((e_i);\mathrm{X};\mathbf{Y})\) or simply \(\mathrm{P}(\mathrm{X};\mathbf{Y})\) , is the minimal polynomial over \(\mathrm{K}(\mathbf{Y})\) of the element \(\sum_{i\in \mathrm{I}}\mathrm{Y}_i e_i\) of \(\mathrm{B}_{(\mathrm{K}(\mathbf{Y}))}\) . *Prove that it belongs to \(\mathrm{K}[\mathrm{X};\mathbf{Y}]\) (cf.~AC, VII, §3, n° 5, p.~230, lemma 1 and théorème 2).
\item
  Let \((e_{i}^{\prime})\) be another basis of B over K and \((\lambda_{ij})\) the matrix with entries in K such that \(e_{i} = \sum_{j} \lambda_{ij} e_{j}^{\prime}\) for \(i \in I\) ; if we set \(Y_{i}^{\prime} = \sum_{j} \lambda_{ji} Y_{j}\) , then \(\mathrm{P}((e_{i}); \mathrm{X}; \mathbf{Y}) = \mathrm{P}((e_{i}^{\prime}); \mathrm{X}; \mathbf{Y}^{\prime})\) . Let \(b = \sum_{i} \beta_{i} e_{i}\) be an element of B; the polynomial \(\mathrm{P}((e_{i}); \mathrm{X}; \beta_{1}, \ldots, \beta_{n})\) is independent of the choice of the basis \((e_{i})\) ; we call it the principal polynomial of b over K. It divides \(\mathrm{Pc}_{\mathrm{B/K}}(b; \mathrm{X})\) and is a multiple of \(\mathrm{Pm}_{\mathrm{B/K}}(b; \mathrm{X}).*\)
\item
  Let \(B'\) be another K-algebra of finite degree; calculate the principal polynomial of the algebra \(B \times B'\) for a suitable basis.
\end{enumerate}

\begin{enumerate}
\def\labelenumi{\arabic{enumi})}
\setcounter{enumi}{2}
\tightlist
\item
  \begin{enumerate}
  \def\labelenumii{\alph{enumii})}
  \tightlist
  \item
    Let A be a central simple algebra of finite degree over K. Prove that the principal polynomial of an element \(a\) of A is \(\mathrm{Pcrd}_{\mathrm{A} / \mathrm{K}}(a; \mathrm{X})\) . The principal polynomial of A (with respect to a basis of A over K) is the polynomial P considered in Lemma 4 of VIII, p.~345 (reduce to the case \(\mathrm{B} = \mathbf{M}_n(\mathrm{K})\) ).
  \end{enumerate}
\end{enumerate}

\begin{enumerate}
\def\labelenumi{\alph{enumi})}
\setcounter{enumi}{1}
\tightlist
\item
  Let B be an absolutely semisimple algebra over K and \(\mathcal{S}\) be the set of classes of simple B-modules. Prove that the principal polynomial of an element b of B is equal to the product of the polynomials \(\mathrm{Pc}_{\mathrm{S}/\mathrm{K}}(b;\mathrm{X})\) for S running through \(\mathcal{S}\) (reduce to the case when K is algebraically closed).
\end{enumerate}

¶ 4) Let D be a field, \(\sigma\) an automorphism of D, A the ring D{[}X{]} \(_{\sigma}\) (VIII, p.~11), \(n\) an integer \(\geqslant 2\) , and \(a\) a nonzero element of D. Denote the element X \(^{n}\) - \(a\) of A by \(f\) . \(a)\) The ideal A \(f\) is two-sided if and only if we have \(\sigma(a) = a\) and \(\sigma^{n}(x) = axa^{-1}\) for every \(x \in \mathrm{D}\) .

\begin{enumerate}
\def\labelenumi{\alph{enumi})}
\setcounter{enumi}{1}
\tightlist
\item
  Suppose that the group \(\mu_{n}\) of n-th roots of unity contained in the center of D is cyclic of order n and fixed by \(\sigma\) . Let \(g \in A\) be an irreducible monic polynomial (VIII, p.~22, Exercise 27) dividing f on the right. Prove that the degree d of g divides n and that f is the product of n/d polynomials of degree d.~(Let h be the monic polynomial such that \(Ah = \bigcap_{\zeta \in \mu_{n}} Ag_{\zeta}\) , where \(g_{\zeta}(X) = g(\zeta X)\) . Observe that we necessarily have \(h(\zeta X) = h(X)\) for every \(\zeta \in \mu_{n}\) , and deduce that h = f.~Then, apply loc. cit., b) by constructing, for \(1 \leqslant i \leqslant n/d\) , an element \(\zeta_{i}\) of \(\mu_{n}\) and a polynomial \(h_{i}\) of degree id such that \(Ag_{\zeta_{1}} \cap \cdots \cap Ag_{\zeta_{i}} = Ah_{i}\) .
\end{enumerate}

If, moreover, the ideal Af is two-sided, then the ring A/Af is semisimple, and its simple components all have the same length (observe that every simple \((\mathrm{A}/\mathrm{Af})\) -module is isomorphic to one of the modules \(A/Ag_{\zeta}\) ).

\begin{enumerate}
\def\labelenumi{\alph{enumi})}
\setcounter{enumi}{2}
\tightlist
\item
  Assume from now on that the hypotheses of b) are satisfied, that the ideal A f is two-sided, and that n is a prime number. Deduce from b) that we are in one of the following situations:
\end{enumerate}

\begin{enumerate}
\def\labelenumi{(\roman{enumi})}
\item
  There exists a \(b \in D\) such that \(a = b^{n}\) and that \(\sigma\) is the inner automorphism associated with b. The ring C = A/Af is isomorphic to the product ring \(D^{n}\) .
\item
  The condition of (i) is not satisfied, but there exists a \(c \in \mathrm{D}\) such that \(a = \sigma^{n-1}(c)\sigma^{n-2}(c)\ldots\sigma(c)c\) . The ring C is simple of length \(n\) .
\item
  There is no element c of D such that \(a = \sigma^{n-1}(c) \sigma^{n-2}(c) \cdots \sigma(c) c\) . The ring C is a field.
\end{enumerate}

In the case when D is commutative, recover the results of Exercise 11 of VIII, p.~330. If, moreover, n = 2, then C is a quaternion algebra over the subfield K of D invariant under \(\sigma\) .

\begin{enumerate}
\def\labelenumi{\alph{enumi})}
\setcounter{enumi}{3}
\item
  Prove that there exists an automorphism \(\tau\) of C of order \(n\) whose subfield of invariants is equal to D. In cases (ii) and (iii) above, D is a Galois subfield of C (Exercise 16 of VIII, p.~274).
\item
  Suppose that \(\sigma\) is the inner automorphism associated with an element \(d\) of D. Prove that we have \(a = d^n z\) , where \(z\) belongs to the center Z of C, and that the ring C is isomorphic to the tensor product \(\mathrm{D} \otimes_{\mathbb{Z}} \mathrm{Z}[\mathrm{X}] / (\mathrm{X}^n - z)\) .
\item
  Suppose that \(\sigma\) is not an inner automorphism. Prove that the center of C is the subfield of Z consisting of the elements invariant under \(\sigma\) .
\end{enumerate}

¶ 5) Let \(K_{0}\) be the field of rational fractions \(\mathbf{Q}(\mathrm{X})\) . Let K be the quadratic extension \(\mathrm{K}_{0}(\rho)\) , where \(\rho\) is a primitive cubic root of unity, D be the cyclic extension \(\mathrm{K}(\sqrt[3]{\mathrm{X}})\) of K, and \(\sigma\) be a generator of the Galois group of D over K. Consider the algebra C defined as in Exercise 4, with n = 3 and a = 2.

\begin{enumerate}
\def\labelenumi{\alph{enumi})}
\item
  Prove that C is a field with center K, of degree 9 over K. (To see that D has no element of norm 2, reduce to proving that there cannot be any relations of the form \(2F^{3} = P^{3} + XQ^{3} + X^{2}R^{3} + 3XPQR\) , where F, P, Q, R are mutually prime polynomials in \(\mathbf{Q}(\rho)[\mathrm{X}]\) . To do this, consider the lowest-degree terms of these polynomials.)
\item
  Prove that there does not exist any automorphism of C extending the unique \(K_{0}\) -automorphism of K distinct from the identity (reduce to the case when such an automorphism preserves D). Deduce that \(K_{0}\) is not a Galois subfield of C (VIII, p.~274, Exercise 16) and that there does not exist a \(K_{0}\) -automorphism of C of order 2.
\item
  Let \(\xi\) be an element of D such that \(\xi^3 = \mathrm{X}\) , and let \(\eta\) be an element of C such that \(\eta^3 = 2\) and \(\eta \xi = \rho \xi \eta\) ; set \(\zeta = \xi + \eta + \eta^2\) . Prove that \(\mathrm{K}(\zeta)\) is a maximal commutative subfield of D and is not Galois over K, whereas the maximal commutative subfields \(\mathrm{K}(\xi)\) and \(\mathrm{K}(\eta)\) are Galois over K (calculate \(\zeta^3\) ).
\end{enumerate}

\begin{enumerate}
\def\labelenumi{\arabic{enumi})}
\setcounter{enumi}{5}
\tightlist
\item
  Let D be a field with center K of degree \(m^{2}\) over K. Let x be an element of D of degree n \textgreater{} 1 over K, and let \(f(\mathrm{X})\) be its minimal polynomial over K.
\end{enumerate}

\begin{enumerate}
\def\labelenumi{\alph{enumi})}
\item
  Prove that a monic polynomial in D{[}X{]} (VIII, p.~11) of degree \textless{} n cannot be right divisible by all polynomials \(X - txt^{-1}\) for \(t \in D^{*}\) (observe that the monic generator of the ideal \(\cap(\mathrm{D}[\mathrm{X}](\mathrm{X}-txt^{-1}))\) has coefficients in K).
\item
  Let \(u(\mathrm{X}), v(\mathrm{X}), w(\mathrm{X})\) be three monic polynomials in D{[}X{]} such that \(u(\mathrm{X}) = v(\mathrm{X})w(\mathrm{X})\) and \(\mathrm{X} - x\) is a right divisor of \(u(\mathrm{X})\) but not of \(w(\mathrm{X})\) . Write \(w(\mathrm{X}) = q(\mathrm{X})(\mathrm{X} - x) + t\) with \(t \in \mathrm{D}^*\) . Prove that \(\mathrm{X} - txt^{-1}\) is a right divisor of \(v(\mathrm{X})\) .
\item
  Deduce from a) and b) that there exist n elements \(t_{1},\ldots,t_{n}\) of D such that we have
\end{enumerate}

\[
f (\mathrm{X}) = \left(\mathrm{X} - t _ {1} x t _ {1} ^ {- 1}\right) \dots \left(\mathrm{X} - t _ {n} x t _ {n} ^ {- 1}\right)
\]

(observe that \(f(\mathrm{X})\) is right divisible by all polynomial \(\mathrm{X} - txt^{-1}\) ).

\begin{enumerate}
\def\labelenumi{\alph{enumi})}
\setcounter{enumi}{3}
\tightlist
\item
  Deduce that \(\operatorname{Trd}_{\mathrm{D}/\mathrm{K}}(x)\) (resp. \(\operatorname{Nrd}_{\mathrm{D}/\mathrm{K}}(x)\) ) is the sum (resp. product) of m elements of the form \(txt^{-1}\) (apply Proposition 8 of VIII, p.~347).
\end{enumerate}

\begin{enumerate}
\def\labelenumi{\arabic{enumi})}
\setcounter{enumi}{6}
\tightlist
\item
  Let D be a field with center K of finite rank over K, endowed with a total order compatible with its additive group structure and such that the product of two positive elements is positive. Prove that D is equal to K. (Prove that K has characteristic zero. If \(D \neq K\) , observe that there exist elements x in D-K such that \(\operatorname{Trd}_{\mathrm{D}/\mathrm{K}}(x) = 0\) , and apply Exercise 6, d)).
\end{enumerate}

\section{Simple Algebras over a Finite Field}

\subsection{Polynomials over a Finite Field}

THEOREM 1 (Chevalley-Warning). --- Let K be a finite commutative field of characteristic p.~Let n be an integer \(\geqslant 1\) and \((f_{i})_{i\in\mathrm{I}}\) a finite family of nonzero elements of \(\mathrm{K}[\mathrm{X}_{1},\ldots,\mathrm{X}_{n}]\) . Denote by Z the set of elements x of \(K^{n}\) such that we have \(f_{i}(\mathbf{x})=0\) for \(i\in I\) . If we have \(n>\sum_{i\in\mathrm{I}}\deg(f_{i})\) , then the cardinal of Z is divisible by p.

Lemma 1. --- Let L be a field, G a finite group, and \(\chi\) a nontrivial homomorphism from G to the multiplicative group \(\mathrm{L}^*\) . We have \(\sum_{x\in \mathrm{G}}\chi (x) = 0\) .

By assumption, there exists an element a of G such that \(\chi(a) \neq 1\) ; since multiplication by a is a permutation of G, we have

\[
\sum_ {x \in \mathrm{G}} \chi (x) = \sum_ {x \in \mathrm{G}} \chi (a x) = \chi (a) \sum_ {x \in \mathrm{G}} \chi (x);
\]

this gives Lemma 1.

Lemma 2. --- Let q be the cardinal of K. For any integer \(m \geqslant 0\) , set \(S_{m} = \sum_{x \in K} x^{m}\) . We have \(S_{m} = -1\) if m is a nonzero multiple of q - 1 and \(S_{m} = 0\) in all other cases.

Recall that \(0^0 = 1\) (I, §2, No.~1, p.~14). Suppose that \(m\) is a multiple of \(q - 1\) . Since the abelian group \(\mathrm{K}^*\) has order \(q - 1\) , we have \(x^m = 1\) for every \(x \in \mathrm{K}^*\) and \(\mathrm{S}_m = 0^m + (q - 1) \cdot 1\) , which gives the assertion in this case.

Suppose that \(m\) is not a multiple of \(q - 1\) . Set \(\chi(x) = x^m\) for \(x \in \mathrm{K}^*\) . Since the multiplicative group \(\mathrm{K}^*\) is cyclic of order \(q - 1\) (V, §12, No.~1, p.~93, Proposition 1), there exists an element \(a\) of \(\mathrm{K}^*\) such that \(\chi(a) \neq 1\) . By Lemma 1 applied to \(\mathrm{G} = \mathrm{K}^*\) , we have

\[
\mathrm{S} _ {m} = 0 ^ {m} + \sum_ {x \in \mathrm{K} ^ {*}} \chi (x) = 0;
\]

this gives Lemma 2.

Let us now prove Theorem 1. Let \(\mathbf{x} = (x_1, \ldots, x_n)\) be an element of \(\mathrm{K}^n\) . We have \(1 - f_i(\mathbf{x})^{q-1} = 0\) if \(f_i(\mathbf{x}) \neq 0\) and \(1 - f_i(\mathbf{x})^{q-1} = 1\) if \(f_i(\mathbf{x}) = 0\) . Set \(\mathrm{P} = \prod_{i=1}^{r}(1 - f_i^{q-1})\) . We have

\[
\mathrm{P} (\mathbf {x}) = \left\{ \begin{array}{l l} 1 & \text { if } \mathbf {x} \in \mathrm{Z}, \\ 0 & \text { if } \mathbf {x} \not \in \mathrm{Z}. \end{array} \right.\tag{1}
\]

Let us expand the polynomial P as \(\sum_{\alpha\in\mathbf{N}^{n}}c_{\alpha}\mathrm{X}^{\alpha}\) ; by assumption, it has degree \(< (q-1)n\) . Let \(\alpha\) be an element of \(N^{n}\) such that \(c_{\alpha}\) is nonzero. Since we have \(\alpha_{1}+\cdots+\alpha_{n}<(q-1)n\) , there exists an integer \(\ell\) such that \(1\leqslant\ell\leqslant n\) and \(0\leqslant\alpha_{\ell}<q-1\) . By Lemma 2, we then have \(\sum_{x\in K}x^{\alpha_{\ell}}=0\) , and therefore

\[
\sum_ {\mathbf {x} \in \mathrm{K} ^ {n}} \mathbf {x} ^ {\alpha} = \prod_ {j = 1} ^ {n} \left(\sum_ {x \in \mathrm{K}} x ^ {\alpha_ {j}}\right) = 0.
\]

We consequently have

\[
\sum_ {\mathbf {x} \in K ^ {n}} \mathrm{P} (\mathbf {x}) = \sum_ {\alpha \in \mathbf {N} ^ {n}} c _ {\alpha} \left(\sum_ {\mathbf {x} \in K ^ {n}} \mathbf {x} ^ {\alpha}\right) = 0.
\]

Now, by formula (1), we have \(\sum_{\mathbf{x}\in\mathrm{K}^{n}}\mathrm{P}(\mathbf{x})=\mathrm{Card}(\mathrm{Z})\cdot1\) , and therefore \(\mathrm{Card}(\mathrm{Z})\cdot1=0\) , which means that \(\mathrm{Card}(\mathrm{Z})\) is divisible by p.

COROLLARY. --- Let V be a vector space of finite dimension n over K and I a finite set, and for each \(i \in I\) , let \(F_{i}: V \to K\) be a homogeneous polynomial mapping of degree \(d_{i} > 0\) . If we have \(\sum_{i \in I} d_{i} < n\) , then there exists a nonzero element x of V such that \(F_{i}(x) = 0\) for every \(i \in I\) .

Let \((e_{1},\ldots,e_{n})\) be a basis of V over K. By the definition of homogeneous polynomial mappings (IV, §5, No.~9, p.~55, Definition 3), for every \(i\in I\) , there exists a homogeneous polynomial \(f_{i}\) in \(\mathrm{K}[\mathrm{X}_{1},\ldots,\mathrm{X}_{n}]\) of degree \(d_{i}\) such that we have \(\mathrm{F}_{i}(\xi_{1}e_{1},\ldots,\xi_{n}e_{n})=f_{i}(\xi_{1},\ldots,\xi_{n})\) . Let Z be the set of elements x of V such that \(\mathrm{F}_{i}(x)=0\) for every \(i\in I\) . By Theorem 1, the cardinal of Z is divisible by p, and since 0 belongs to S, we have \(\mathrm{Card}(\mathrm{Z})\geqslant p>1\) .

\subsection{Simple Algebras over Finite Fields}

THEOREM 2 (Wedderburn). --- Every finite field is commutative.

Let D be a finite field, and let K be its center. The K-algebra D is central simple of degree \(m^{2}\) , where m is a strictly positive integer. The reduced norm is a homogeneous polynomial mapping \(Nrd : D \to K\) of degree m (VIII, p.~345, Proposition 6), and we have \(\text{Nrd}(a) \neq 0\) for every \(a \neq 0\) in D (VIII, p.~340, Proposition 3). The corollary above implies that \(m \geqslant m^{2}\) , and therefore m = 1. So we have D = K.

COROLLARY 1. --- Every finite simple ring is isomorphic to a matrix ring \(\mathbf{M}_{n}(\mathrm{L})\) , where n is a strictly positive integer and L is a finite commutative field.

This follows from Theorem 2 and the structure theorem for simple rings (VIII, p.~120, Theorem 1).

COROLLARY 2. --- Let K be a finite commutative field. Every central simple algebra over K is isomorphic to a matrix algebra \(\mathbf{M}_{n}(\mathrm{K})\) , where n is a strictly positive integer.

This follows from Theorem 2 and the structure theorem for central simple algebras (VIII, p.~252, Theorem 1).

Remarks. --- 1) Here is another proof of Theorem 2. Let D be a finite field, let K be its center, and let L be a maximal commutative subfield of D. Let x be an element of \(D^{*}\) . It belongs to a maximal commutative subfield \(L_{1}\) of D. We have the equality

\[
[ \mathrm{D}: \mathrm{K} ] = [ \mathrm{L}: \mathrm{K} ] ^ {2} = [ \mathrm{L} _ {1}: \mathrm{K} ] ^ {2}
\]

by Corollary 2 of VIII, p.~265, and therefore \([\mathrm{L}:\mathrm{K}] = [\mathrm{L}_1:\mathrm{K}]\) . By Proposition 3 of V, §12, No.~2, p.~94, the extensions L and \(\mathrm{L}_1\) of K are isomorphic. By VIII, p.~263, Corollary, there exists an element \(a\) of \(\mathrm{D}^*\) such that \(a\mathrm{La}^{-1} = \mathrm{L}_1\) , so \(a^{-1}xa\) belongs to L. We then have \((ay)^{-1}x(ay) = a^{-1}xa\) for every \(y\in \mathrm{L}^*\) . Consequently, if S is a set of representatives of the left cosets of \(\mathrm{D}^*\) modulo \(\mathrm{L}^*\) , then every element of \(\mathrm{D}^* -\{1\}\) can be written as \(sxs^{-1}\) with \(s\in \mathrm{S}\) and \(x\in \mathrm{L}^* -\{1\}\) . We denote the order of \(\mathrm{D}^*\) by \(d\) and that of \(\mathrm{L}^*\) by \(\ell\) . Since the cardinal of S is equal to \(d / \ell\) , we have \(d - 1\leqslant (d / \ell)(\ell -1) = d - d / \ell\) . It follows that \(\ell = d\) , and therefore \(\mathrm{L} = \mathrm{D}\) , which proves that the field D is commutative.

\begin{enumerate}
\def\labelenumi{\arabic{enumi})}
\setcounter{enumi}{1}
\tightlist
\item
  Let L be a commutative field with the following property:
\end{enumerate}

(C \(_{1}\) ) Let V be a vector space of finite dimension n over the field L, and let d be such that 0 \textless{} d \textless{} n.~For every homogeneous polynomial mapping F : V → L of degree d, there exists a nonzero element x of V such that F(x) = 0.

The proof of Theorem 2 shows that every field of finite degree over L with center L is equal to L.

By the corollary of VIII, p.~356, every finite field has property \((\mathrm{C}_{1})\) . We can prove (VIII, p.~360, Exercise 7) that the following fields have property \((\mathrm{C}_{1})\) :

- every algebraic extension of a finite field

-- every field of rational fractions in one variable with coefficients in an algebraically closed field (Tsen's theorem).

- *every field endowed with a discrete valuation for which it is complete and whose residue field is algebraically closed (VIII, p.~332, Exercise 17).*

\begin{enumerate}
\def\labelenumi{\arabic{enumi})}
\setcounter{enumi}{2}
\tightlist
\item
  Suppose that the field \(\mathbf{K}\) satisfies the following condition:
\end{enumerate}

- If \(\mathrm{L}\) is an extension of \(\mathrm{K}\) of finite degree, then it is cyclic and the norm mapping \(\mathrm{N}:\mathrm{L}^{*}\to \mathrm{K}^{*}\) is surjective.

This condition is satisfied, in particular, when the field K is finite (V, §12, No.~2, p.~95, Proposition 4). We can then prove that every field of finite degree over K with center K is equal to K (Exercise 10 of VIII, p.~329).

\BourbakiExercises{Exercises}

\begin{enumerate}
\def\labelenumi{\arabic{enumi})}
\tightlist
\item
  Let \(D\) be a field.
\end{enumerate}

\begin{enumerate}
\def\labelenumi{\alph{enumi})}
\item
  Let E be a proper subfield of D whose multiplicative group \(E^{*}\) has finite index in \(D^{*}\) . Prove that D is finite (for a sequence \((a_{n})\) of distinct elements of E and an element x of D---E, consider the classes mod \(E^{*}\) of the elements \(x + a_{n}\) ).
\item
  Prove that an element of D that has only finitely many conjugates belongs to the center Z of D (apply a) to the commutant of this element in D, which is a field).
\item
  Deduce that a polynomial in Z{[}X{]} that has a root in D-Z has infinitely many roots in D-Z.
\end{enumerate}

¶ 2) Let A be a K-algebra and m an integer. Suppose that every element of A is algebraic of degree \(\leqslant m\) over K.

\begin{enumerate}
\def\labelenumi{\alph{enumi})}
\item
  Prove that if A is a primitive ring, then there exist an integer \(r \leqslant m\) and a field D whose elements all have degree \(\leqslant m\) over K such that A is isomorphic to the matrix algebra \(\mathbf{M}_{r}(\mathrm{D})\) . (If A is not simple, then deduce from the density theorem that for every integer n, there exist a subalgebra \(A_{n}\) of A and a surjective homomorphism \(A_{n} \to M_{n}(\mathrm{D})\) . Observe that the algebra \(\mathbf{M}_{n}(\mathrm{D})\) contains elements X such that \(X^{n-1} \neq 0\) and \(X^{n} = 0\) .)
\item
  Suppose that the field K is infinite. Prove that the algebra \(A/\Re(A)\) is semisimple and admits at most m simple components (observe that the algebra \(K^{n}\) contains elements of degree n over K; deduce that A admits at most m maximal two-sided ideals, and conclude using a)). Prove that if K is, moreover, perfect, then \(A/\Re(A)\) has finite rank over K.
\end{enumerate}

¶ 3) Let K be a commutative field. We say that a K-algebra A is algebraic if all its elements are algebraic over K. Every algebra of finite rank is algebraic; every subalgebra and quotient algebra of an algebraic algebra is algebraic. The radical of an algebraic K-algebra is a nil ideal (VIII, p.~166, Exercise 5).

\begin{enumerate}
\def\labelenumi{\alph{enumi})}
\item
  Prove that an algebraic algebra is a pseudoregular ring (VIII, p.~178, Exercise 4; observe that for every \(x \in \mathrm{A}\) , there exist an element \(y \in \mathrm{K}[x]\) and an integer \(k\) such that \(x^k = x^{k+1}y\) ).
\item
  Let A be an algebraic K-algebra without any nilpotent elements other than 0. Prove that A is isomorphic to a subalgebra of a product of fields. (Observe that every quotient algebra of A has the same properties. Prove that if A is primitive, then it is a field, by reasoning as in Exercise 2, a). In the general case, use Exercise 12 of VIII, p.~168.)
\end{enumerate}

¶ 4) Let K be a finite field and A an algebraic K-algebra (Exercise 3).

\begin{enumerate}
\def\labelenumi{\alph{enumi})}
\item
  Suppose that A is a field. Prove that A is commutative. (Let \(x\) be a noncentral element of A. Consider the field \(Z(x)\) , where \(Z\) is the center of A, and prove that there exists an element \(y \in A\) such that \(yxy^{-1} = x^r\) , with \(x^r \neq x\) . Deduce a contradiction by considering the subfield \(K(x,y)\) of A.)
\item
  Suppose that A does not contain any nonzero nilpotent element. Prove that the algebra A is commutative (use a) and Exercise 3).
\end{enumerate}

\begin{enumerate}
\def\labelenumi{\arabic{enumi})}
\setcounter{enumi}{4}
\item
  Let K be a finite field, m an integer, and A a K-algebra without radical whose elements all have degree \(\leqslant m\) . Prove that there exists a finite extension L of K such that A is isomorphic to a subalgebra of a product \((\mathbf{M}_{r}(\mathrm{L}))^{\mathrm{I}}\) whose projections are simple algebras (cf.~Exercise 2, a)). Also prove the converse.
\item
  Let A be a ring such that for every \(a \in \mathrm{A}\) , there exists an integer \(n(a) > 1\) such that \(a^{n(a)} = a\) . Prove that A is isomorphic to a product of finitely many reduced commutative algebraic algebras over finite fields. (Observe that A is annihilated by an integer \(n\) , and then that for every prime number \(p\) , the \(p\) -component \(\mathrm{A}_p\) of the additive group A is a two-sided ideal annihilated by \(p\) . Apply Exercise 4, b).) Also prove the converse.
\item
  \begin{enumerate}
  \def\labelenumii{\alph{enumii})}
  \tightlist
  \item
    Prove that if the field K has property \((\mathrm{C}_{1})\) (VIII, p.~357), then so does every algebraic extension of K. (Reduce to the case of an extension L of K of finite degree r. Prove that if F is a homogeneous polynomial function of degree d on an L-vector space V, then the function \(N_{L/K} \circ F\) is homogeneous polynomial of degree rd on the K-vector space V.)
  \end{enumerate}
\end{enumerate}

\(^{*}b)\) Assume from now on that K is algebraically closed. Let \(F_{1},\ldots,F_{n}\) be homogeneous polynomials of degree \textgreater0 in m variables. Prove that if n \textless{} m, then there exists an \(a \neq 0\) in \(K^{n}\) such that \(F_{1}(a) = \cdots = F_{n}(a) = 0\) (use AC, VIII, §2, n° 4, p.~19, corollaire 1 du théorème 3 and AC, VIII, §3, n° 1, p.~24, proposition 2). c) Deduce from b) that the field \(K(X)\) has property \((C_{1})\) .

(Let f be a homogeneous polynomial of degree d in n variables whose coefficients are polynomial in X, and let \(\nu\) be the maximum of the degrees of these coefficients. Let N be an integer, and let \(G_{1},\ldots,G_{n}\) be polynomials in X of degree \(\leqslant N\) . Prove that the equality \(\mathrm{F}(\mathrm{G}_{1}(\mathrm{X}),\ldots,\mathrm{G}_{n}(\mathrm{X}))=0\) is equivalent to the annihilation of \(Nd+\nu+1\) homogeneous polynomials of degree \textgreater0 in the \(n(N+1)\) coefficients of the \(G_{i}\) .)

\(d\) ) Deduce from \(a\) ) and \(c\) ) that an extension of an algebraically closed field of transcendence degree \(\leqslant 1\) has property \((\mathrm{C}_1)\) (``Tsen's theorem'').*

\section{Quaternion Algebras}

In this section, K is a commutative field.

\subsection{General Properties of Quaternion Algebras}

Let \(\alpha, \beta, \gamma\) be elements of K, and let F be the quaternion algebra of type \((\alpha, \beta, \gamma)\) . Recall (III, §2, No.~5, p.~445) \(^{(1)}\) that F is an associative unital K-algebra that has a basis \((1, i, j, k)\) over K satisfying the relations

\[
i ^ {2} = \alpha + \beta i, j ^ {2} = \gamma , i j = k, j i = \beta j - k.\tag{1}
\]

It is a Cayley algebra (III, §2, No.~4, p.~441, Definition 1) whose conjugation satisfies

\[
\overline {{i}} = \beta - i, \quad \overline {{j}} = - j, \quad \overline {{k}} = - k.\tag{2}
\]

Recall that the Cayley trace and norm of F are the mappings \(T_{F}\) and \(N_{F}\) from F to K defined by \(\mathrm{T}_{\mathrm{F}}(q)=q+\overline{q}\) and \(\mathrm{N}_{\mathrm{F}}(q)=q\overline{q}\) .

The linear subspace E of F with basis \((1,i)\) is a commutative Cayley subalgebra of F; it is a quadratic algebra of type \((\alpha,\beta)\) , and F can be identified with the Cayley extension of E defined by \(\gamma\) (III, §2, No.~5, p.~444). For every \(z\in E\) , we have \(zj=j\overline{z}\) . Every element q of F can be written uniquely as \(x+jy\) with \(x,y\in E\) , and we have

\[
\overline {{q}} = \overline {{x}} - j y, \quad \mathrm{T} _ {\mathrm{F}} (q) = x + \overline {{x}}, \quad \mathrm{N} _ {\mathrm{F}} (q) = x \overline {{x}} - \gamma y \overline {{y}}.\tag{3}
\]

PROPOSITION 1. --- The characteristic polynomial of an element q of F is equal to \(\left(\mathrm{X}^{2}-\mathrm{T}_{\mathrm{F}}(q)\mathrm{X}+\mathrm{N}_{\mathrm{F}}(q)\right)^{2}\) .

By the above, the algebra \(\mathrm{F}\) is a free right \(\mathrm{E}\) -module with basis \((1,j)\) . Consequently, \(\mathrm{F}[\mathrm{X}]\) is a free right \(\mathrm{E}[\mathrm{X}]\) -module with basis \((1,j)\) . We denote by \(u\) the endomorphism of the right \(\mathrm{E}[\mathrm{X}]\) -module \(\mathrm{F}[\mathrm{X}]\) defined by \(u(\mathrm{P}) = (\mathrm{X} - q)\mathrm{P}\) for \(\mathrm{P} \in \mathrm{F}[\mathrm{X}]\) . The characteristic polynomial of \(q\) is the determinant of \(u\) viewed as an endomorphism of the \(\mathrm{K}[\mathrm{X}]\) -module \(\mathrm{F}[\mathrm{X}]\) . By Proposition 6 of III, §9, No.~4, p.~546, it is equal to \(\mathrm{N}(\det u)\) , where \(\mathrm{N}\) denotes the norm from \(\mathrm{E}[\mathrm{X}]\) to \(\mathrm{K}[\mathrm{X}]\) . Let us write \(q\) as \(x + jy\) with \(x, y \in \mathrm{E}\) . The matrix of \(u\) with respect to the basis \((1,j)\) is \(\begin{pmatrix} \mathrm{X} - x & -\gamma \overline{y} \\ -y & \mathrm{X} - \overline{x} \end{pmatrix}\) ; its determinant is equal to \(\mathrm{D} = (\mathrm{X} - x)(\mathrm{X} - \overline{x}) - \gamma y\overline{y} = \mathrm{X}^2 - \mathrm{T}_{\mathrm{F}}(q)\mathrm{X} + \mathrm{N}_{\mathrm{F}}(q)\) (cf.~formula (3)). Since \(\mathrm{D}\) belongs to \(\mathrm{K}[\mathrm{X}]\) , we have \(\mathrm{N}(\mathrm{D}) = \mathrm{D}^2\) ; Proposition 1 follows.

Remarks. --- 1) Suppose that the characteristic of K is different from 2, and set \(i' = 2i - \beta\) . Then \((1, i')\) is a basis of E over K, and we have \(i'^2 = 4\alpha + \beta^2\) . It follows that E is isomorphic to the quadratic algebra of type \((4\alpha + \beta^2, 0)\) and F to the quaternion algebra of type \((4\alpha + \beta^2, 0, \gamma)\) .

\begin{enumerate}
\def\labelenumi{\arabic{enumi})}
\setcounter{enumi}{1}
\item
  The quaternion algebra of type \((\alpha,0,\gamma)\) is isomorphic to the quaternion algebra of type \((\gamma,0,\alpha)\) (III, §2, No.~5, p.~445). It is also isomorphic to the quaternion algebra of type \((\alpha a^{2},0,\gamma c^{2})\) for every pair \((a,c)\) of nonzero elements of K.
\item
  Let q be an element of F. Then q is nilpotent if and only if its characteristic polynomial is equal to \(X^{4}\) , that is, \(\mathrm{T}_{\mathrm{F}}(q)=\mathrm{N}_{\mathrm{F}}(q)=0\) ; we then have \(q^{2}=0\) .
\end{enumerate}

Example. --- The matrix algebra \(\mathbf{M}_{2}(\mathrm{K})\) is isomorphic to the quaternion algebra of type \((0,1,1)\) . Indeed, consider the quadratic algebra \(E = K \times K\) (of type \((0,1)\) ) and the quaternion algebra \(F = E + Ej\) , which is the Cayley extension of E defined by the element \(\gamma = 1\) . The mapping \((a,b) \mapsto \begin{pmatrix} a & 0 \\ 0 & b \end{pmatrix}\) is an algebra homomorphism from E to \(\mathbf{M}_{2}(\mathrm{K})\) . Since for a, b in K, we have

\[
\left( \begin{array}{c c} 0 & 1 \\ 1 & 0 \end{array} \right) \left( \begin{array}{c c} 0 & 1 \\ 1 & 0 \end{array} \right) = \left( \begin{array}{c c} 1 & 0 \\ 0 & 1 \end{array} \right), \quad \left( \begin{array}{c c} 0 & 1 \\ 1 & 0 \end{array} \right) \left( \begin{array}{c c} a & 0 \\ 0 & b \end{array} \right) = \left( \begin{array}{c c} b & 0 \\ 0 & a \end{array} \right) \left( \begin{array}{c c} 0 & 1 \\ 1 & 0 \end{array} \right),
\]

this homomorphism extends to an algebra homomorphism \(\theta : \mathrm{F} \longrightarrow \mathbf{M}_2(\mathrm{K})\) defined by

\[
\theta \bigl ((a, b) + (c, d) j \bigr) = \left( \begin{array}{c c} a & c \\ d & b \end{array} \right)  .
\]

This homomorphism is bijective. When the characteristic of \(\mathbf{K}\) is different from 2, the algebra \(\mathbf{M}_2(\mathbf{K})\) is also isomorphic to the quaternion algebra of type (1,0,1) (Remark 1).

\subsection{The Center of a Quaternion Algebra}

Let \(\alpha, \beta, \gamma\) be elements of \(K\) , and let \(F\) be the quaternion algebra of type \((\alpha, \beta, \gamma)\) .

PROPOSITION 2. --- a) Suppose that the field K has characteristic different from 2. If \(\gamma\) or \(4\alpha + \beta^{2}\) is nonzero, then the center of F is equal to K; otherwise, it has dimension 2 and is generated by 1 and ij - ji.

\begin{enumerate}
\def\labelenumi{\alph{enumi})}
\setcounter{enumi}{1}
\tightlist
\item
  Suppose that the field \(\mathrm{K}\) has characteristic 2. If \(\beta \neq 0\) , then the center of \(\mathrm{F}\) is equal to \(\mathrm{K}\) ; if \(\beta = 0\) , then the algebra \(\mathrm{F}\) is commutative.
\end{enumerate}

By formula (30) of III, §2, No.~5, p.~444, we have

\[
i j - j i = - \beta j + 2 k, \quad j k - k j = \beta \gamma - 2 \gamma i, \quad k i - i k = - 2 \alpha j - \beta k. \tag {4}
\]

An element \(q = x + yi + zj + tk\) of F is central if and only if it commutes with i and j, that is, we have

\[
2 z + \beta t = - \beta z + 2 \alpha t = 0 \quad \text { and } \quad 2 \gamma t = \beta \gamma t = 2 y = \beta y = 0.\tag{5}
\]

First suppose that the characteristic of K is different from 2. If \(\gamma\) is nonzero, then the equalities in (5) imply y = t = 0 and then z = 0; consequently, we have \(q \in K\) . If \(\gamma = 0\) and \(4\alpha + \beta^{2} \neq 0\) , then they imply y = z = t = 0, so we have \(q \in K\) . If \(\gamma = 4\alpha + \beta^{2} = 0\) , then the system (5) reduces to \(y = 2z + \beta t = 0\) , so we have \(q = x + t/2(ij - ji)\) , which completes the proof of a).

Now suppose that the field K has characteristic 2. The system (5) can then be written as \(\beta t = \beta z = \beta y = 0\) ; assertion b) follows.

\subsection{Simplicity of Quaternion Algebras}

PROPOSITION 3. --- Let \(\alpha, \beta, \gamma\) be elements of K, and let \(f\) be a quaternion algebra of type \((\alpha, \beta, \gamma)\) . Denote its Cayley trace and norm by \(\mathrm{T}_{\mathrm{F}}\) and \(\mathrm{N}_{\mathrm{F}}\) . The following properties are equivalent:

\begin{enumerate}
\def\labelenumi{(\roman{enumi})}
\item
  The algebra \(\mathbf{F}\) is central simple.
\item
  For every nonzero element \(x\) of \(\mathrm{F}\) , there exists an element \(y\) in \(\mathrm{F}\) such that \(\mathrm{T}_{\mathrm{F}}(xy) \neq 0\) .
\item
  We have \((4\alpha + \beta^{2})\gamma \neq 0\) .
\end{enumerate}

Suppose that these properties hold. Then for every x in F, the reduced characteristic polynomial of x is \(\mathrm{X}^{2}-\mathrm{T}_{\mathrm{F}}(x)\mathrm{X}+\mathrm{N}_{\mathrm{F}}(x)\) . In particular, \(\mathrm{T}_{\mathrm{F}}(x)\) is the reduced trace of x, and \(\mathrm{N}_{\mathrm{F}}(x)\) is its reduced norm.

(i)⇒(ii): If the algebra F is central simple, then it follows from Proposition 1 of VIII, p.~361 and the definition of the reduced trace (VIII, p.~340, Definition 2) that \(T_{F}\) is its reduced trace; assertion (ii) follows (VIII, p.~343, Proposition 5).

\begin{enumerate}
\def\labelenumi{(\roman{enumi})}
\setcounter{enumi}{1}
\tightlist
\item
  \(\Leftrightarrow\) (iii): Let \((e_i)_{1\leqslant i\leqslant 4}\) be a basis of F of type \((\alpha ,\beta ,\gamma)\) (III, §2, No.~5, p.~445). The matrix \((\mathrm{T_F}(e_i e_j))\) is equal to
\end{enumerate}

\[
\left( \begin{array}{c c c c} 2 & \beta & 0 & 0 \\ \beta & 2 \alpha + \beta^ {2} & 0 & 0 \\ 0 & 0 & 2 \gamma & \beta \gamma \\ 0 & 0 & \beta \gamma & - 2 \alpha \gamma \end{array} \right).
\]

Its determinant is \(-\gamma^{2}(4\alpha + \beta^{2})^{2}\) . The equivalence of properties (ii) and (iii) follows from V, §8, No.~2, p.~49, Lemma 1.

(iii)⇒(i): Suppose \((4\alpha + \beta^{2})\gamma \neq 0\) . We then have \(\gamma \neq 0\) , and we have \(\beta \neq 0\) if K has characteristic 2. By Proposition 2, the algebra F is central. Let x be an element of the Jacobson radical of F. For every \(y \in F\) , xy is nilpotent, so \(T_{F}(xy) = 0\) (Remark 3 of VIII, p.~362). Since (ii) is equivalent to (iii), we have x = 0. This proves that F is a semisimple K-algebra. Since its center is K, it is simple.

The last assertion follows from Proposition 1 of VIII, p.~361 and the definition of the reduced characteristic polynomial (VIII, p.~340, Definition 1).

Denote the characteristic of K by p.~By Proposition 3, if \(p \neq 2\) , then every quaternion algebra over K of type \((\alpha, 0, \gamma)\) with \(\alpha\) and \(\gamma\) in \(K^{*}\) is central simple. If p = 2, then every quaternion algebra of type \((\alpha, 1, \gamma)\) with \(\alpha \in K\) and \(\gamma \in K^{*}\) is central simple. Conversely, we have the following.

PROPOSITION 4. --- Let A be a central simple algebra of degree 4 over K. Denote the characteristic of K by p.

\begin{enumerate}
\def\labelenumi{\alph{enumi})}
\item
  If \(p \neq 2\) , then there exist nonzero elements \(\alpha\) and \(\gamma\) of \(K\) such that the algebra \(A\) is isomorphic to the quaternion algebra of type \((\alpha, 0, \gamma)\) .
\item
  If \(p = 2\) , then there exist an element \(\alpha\) of \(K\) and an element \(\gamma\) of \(K^*\) such that the algebra \(A\) is isomorphic to the quaternion algebra of type \((\alpha, 1, \gamma)\) .
\end{enumerate}

By Wedderburn's theorem (VIII, p.~120, Theorem 1), there exist an integer \(r \geqslant 1\) and a field D with center K such that A is isomorphic to \(\mathbf{M}_r(\mathrm{D})\) . We then have \(r^2[\mathrm{D}:\mathrm{K}] = [\mathrm{A}:\mathrm{K}] = 4\) . If \(r = 2\) , then A is isomorphic to \(\mathbf{M}_2(\mathrm{K})\) , and Proposition 4 follows from the example of VIII, p.~362. Otherwise, we have \(r = 1\) , and A is a field with center K. It then has a maximal commutative subfield E that is a separable extension of K; since

A has degree 4 over K, the extension E has degree 2 over K (VIII, p.~265, Corollary 2). It is therefore quadratic (III, §2, No.~3, p.~439). Let s be the conjugation of E (III, §2, No.~3, p.~440). By the Skolem--Noether theorem (VIII, p.~256, Corollary 1), there exists an invertible element j of A such that we have \(s(x) = jxj^{-1}\) for every x in E. The field E is separable over K, so we have \(s \neq Id_{E}\) , so that \(j \notin E\) . Since A is a vector space of dimension 4 over K, it is a left vector space of dimension 2 over E, so we have \(A = E \oplus Ej\) . We have \(s^{2} = Id_{E}\) , so the element \(j^{2}\) of A belongs to the center of A; hence there exists an element \(\gamma\) of \(K^{*}\) such that \(j^{2} = \gamma\) .

When \(p \neq 2\) , there exist an element i of E and an element \(\alpha \in K^{*}\) such that \(\mathrm{E} = \mathrm{K}(i)\) and \(i^{2} = \alpha\) (V, §11, No.~9, p.~93, Example 3); in this case, A is isomorphic to the quaternion algebra of type \((\alpha, 0, \gamma)\) . When p = 2, there exist an element i of E and an element \(\alpha\) of K such that \(\mathrm{E} = \mathrm{K}(i)\) and \(i^{2} = i + \alpha\) (V, §11, No.~9, p.~93, Example 2), so that A is isomorphic to the quaternion algebra of type \((\alpha, 1, \gamma)\) .

COROLLARY 1. --- Let A be a central simple K-algebra of finite degree \textgreater{} 1 whose elements are all algebraic of degree \(\leqslant 2\) over K. Then A is isomorphic to a quaternion algebra over K.

If K is finite, then the algebra A is isomorphic to a matrix algebra \(\mathbf{M}_{n}(\mathrm{K})\) (VIII, p.~357, Corollary 2) and therefore contains elements of degree n over K; the assumption implies n = 2 and therefore the result in this case (VIII, p.~362, Example). Suppose that the field K is infinite. Let L be a maximal étale subalgebra of A. By V, §7, No.~4, p.~41, Proposition 7, there exists an element x of A such that the K-algebra L is equal to \(K[x]\) , hence by assumption has degree \(\leqslant 2\) . Since we have \([A : K] = [L : K]^{2}\) (VIII, p.~264, Proposition 4 and p.~262, Proposition 3), we conclude that \([A : K] = 4\) . Corollary 1 then follows from Proposition 4.

COROLLARY 2. --- Let \((\mathrm{E}, s)\) be a Cayley algebra over K such that the K-algebra E is central simple of finite degree \textgreater{} 1 over K. Then E is isomorphic to a quaternion algebra over K.

Every element \(u\) of \(\mathrm{E}\) satisfies \(u^2 - \mathrm{T}_{\mathrm{E}}(u)u + \mathrm{N}_{\mathrm{E}}(u) = 0\) , so the K-algebra \(\mathrm{E}\) is isomorphic to a quaternion algebra (Corollary 1).

\subsection{Criteria for a Quaternion Algebra to Be a Field}

Let \(\alpha,\beta,\gamma\) be elements of the field K, and let F be the quaternion algebra of type \((\alpha,\beta,\gamma)\) . As in No.~1, we denote the canonical basis of F by \((1,i,j,k)\) and the subalgebra \(K+Ki\) of F by E.

PROPOSITION 5. --- The following properties are equivalent:

\begin{enumerate}
\def\labelenumi{(\roman{enumi})}
\item
  The quaternion algebra \(\mathbf{F}\) is a field.
\item
  There is no nonzero element \(q \in \mathrm{F}\) such that \(\mathrm{T}_{\mathrm{F}}(q) = \mathrm{N}_{\mathrm{F}}(q) = 0\) .
\item
  Every element of F of square zero is zero.
\item
  There does not exist any nonzero vector \((x,y,z,t)\) in \(K^{4}\) such that
\end{enumerate}

\[
x ^ {2} + \beta x y - \alpha y ^ {2} - \gamma (z ^ {2} + \beta z t - \alpha t ^ {2}) = 0.
\]

\begin{enumerate}
\def\labelenumi{(\alph{enumi})}
\setcounter{enumi}{21}
\tightlist
\item
  There does not exist any nonzero vector \((x,y,z)\) in \(K^{3}\) such that
\end{enumerate}

\[
x ^ {2} + \beta x y - \alpha y ^ {2} - \gamma z ^ {2} = 0.
\]

\begin{enumerate}
\def\labelenumi{(\roman{enumi})}
\setcounter{enumi}{5}
\tightlist
\item
  The quadratic algebra E is a field, and \(\gamma\) is not the norm of an element of E.
\end{enumerate}

An element q of F is invertible if and only if \(\mathrm{N}_{\mathrm{F}}(q)\) is not 0. The equivalence of (i) and (iv) therefore follows from formula (31) of III, §2, No.~5, p.~445; it is clear that (i) implies (iii) and that (iv) implies (v).

The equivalence of (ii) and (iii) follows from Remark 3 of VIII, p.~362.

Suppose that F is not a field. If \(\gamma(4\alpha+\beta^{2})\neq0\) , then the algebra F is central simple of degree 4 over K; it is isomorphic to the algebra \(\mathbf{M}_{2}(\mathrm{K})\) (VIII, p.~120, Theorem 1) and therefore contains a nonzero element of square zero. If \(\gamma=0\) , then we have \(j^{2}=0\) . If \(4\alpha+\beta^{2}=0\) , then we have \((2i-\beta)^{2}=0\) and \(2i-\beta\neq0\) if K has characteristic different from 2. Finally, if K has characteristic 2 and \(\beta\) is zero, then we have \(\mathrm{T}_{\mathrm{F}}(q)=0\) for every \(q\in F\) (III, §2, No.~5, p.~445, formula (31)). Since F is not a field, there exists a nonzero element q of F such that \(\mathrm{N}_{\mathrm{F}}(q)=0\) ; we have \(q^{2}=0\) . This proves the implication \((\mathrm{iii})\Rightarrow(\mathrm{i})\) .

Let \(q = x + yi\) be an element of E. We have \(\mathrm{N}_{\mathrm{E / K}}(q) = x^2 + \beta xy - \alpha y^2\) . Suppose that property (v) holds. We have \(\mathrm{N}_{\mathrm{E / K}}(q) - \gamma \neq 0\) and \(\mathrm{N}_{\mathrm{E / K}}(q) \neq 0\) if \(q \neq 0\) , and therefore (vi).

Finally, suppose that property (vi) holds. Let q be a nonzero element of F; we write it as \(u + vj\) with u and v in E. If v is zero, then q is invertible. If v is not zero, then we have \(\mathrm{N}_{\mathrm{F}}(q) = \mathrm{N}_{\mathrm{F}}(v)\mathrm{N}_{\mathrm{F}}(v^{-1}u + j) = \mathrm{N}_{\mathrm{E}/\mathrm{K}}(v)\big(\mathrm{N}_{\mathrm{E}/\mathrm{K}}(v^{-1}u) - \gamma\big)\) by III, §2, No.~5, p.~443, formula (24). Since \(\gamma\) is not a norm, \(\mathrm{N}_{\mathrm{F}}(q)\) is not zero, and q is invertible.

Remark. --- Suppose that the quaternion algebra F is a field. It follows from the equality \(j^{2} = \gamma\) that we have \(\gamma \neq 0\) . By Proposition 2 of VIII, p.~363, the center of F is equal to K unless K has characteristic 2 and \(\beta\) is zero, in which case the algebra F is commutative.

\subsection{Algebras over Maximal Ordered Fields}

Let R be a maximal ordered field (VI, §2, No.~5, p.~25). Let C be the quadratic R-algebra of type \((-1,0)\) ; if \((1,i)\) is its canonical basis, then we have \(i^{2}=-1\) . Moreover, C is an algebraic closure of R (VI, §2, No.~6, p.~26, Theorem 3). Let H be the quaternion R-algebra of type \((-1,0,-1)\) . The multiplication table of H in its canonical basis \((1,i,j,k)\) is given by

\[
i ^ {2} = j ^ {2} = k ^ {2} = - 1, \quad i j = - j i = k, \quad - i k = k i = j, \quad j k = - k j = i.
\]

We identify C with the subalgebra \(R + Ri\) of H. The conjugate of an element \(q = x + yi + zj + tk\) of H is \(\overline{q} = x - yi - zj - tk\) . The Cayley trace and norm of q are given by

\[
\mathrm{T} (q) = q + \overline {{q}} = 2 x, \quad \mathrm{N} (q) = q \overline {{q}} = x ^ {2} + y ^ {2} + z ^ {2} + t ^ {2}.
\]

Since R is an ordered field, we have \(N(q) > 0\) if \(q \neq 0\) , so H is a field, with center R (VIII, p.~363, Proposition 2). The reduced trace and norm of an element q of H are \(T(q)\) and \(N(q)\) , respectively.

THEOREM 1. --- Let D be an R-algebra of finite degree that is a field. Then D is isomorphic to R, C, or H.

Denote the center of D by Z, and let L be a maximal commutative subfield of D. We have \([D:Z]=[L:Z]^{2}\) by VIII, p.~265, Corollary 2; we also have \([L:R]\leqslant2\) because C is an algebraic closure of R. There are consequently three possible cases:

\begin{enumerate}
\def\labelenumi{\alph{enumi})}
\item
  We have \(\mathrm{R} = \mathrm{Z} = \mathrm{L}\) , so \([\mathrm{D}:\mathrm{Z}] = 1\) and \(\mathrm{D} = \mathrm{R}\) .
\item
  We have \(\mathrm{R} \neq \mathrm{Z}\) and \(\mathrm{Z} = \mathrm{L}\) , so \([\mathrm{D} : \mathrm{Z}] = 1\) and \(\mathrm{D} = \mathrm{L}\) . In this case, \(\mathrm{D}\) is isomorphic to \(\mathrm{C}\) .
\item
  We have R = Z and \([L : R] = 2\) , so \([D : R] = 4\) . By Proposition 4 of VIII, p.~364, the R-algebra D is isomorphic to a quaternion algebra of type \((\alpha, 0, \gamma)\) , where \(\alpha\) and \(\gamma\) are nonzero elements of R. Let \(i \in D - Z\) be such that \(i^{2} = \alpha\) . We have \(\alpha \neq 0\) . If \(\alpha > 0\) , then there exists an \(a \in R\) such that \(a^{2} = \alpha\) (VI, §2, No.~6, p.~26, Theorem 3); we then have \((a - i)(a + i) = 0\) , which is absurd because D is a field. So we have \(\alpha < 0\) . The inequality \(\gamma < 0\) is shown analogously. There then exist elements a and c of \(R^{*}\) such that \(\alpha = -a^{2}\) and \(\gamma = -c^{2}\) (loc. cit.). The algebra D is therefore isomorphic to the quaternion algebra of type \((-1,0,-1)\) (VIII, p.~362, Remark 2), that is, to H.
\end{enumerate}

Remarks. --- 1) Let O be the octonion algebra of type \((-1,0,-1,-1)\) over R (III, Appendix, No.~3, p.~615). Let D be an alternative Cayley algebra over R such that every nonzero element of D has an inverse. We can prove (VIII, p.~370, Exercise 5) that D is isomorphic to R, C, H, or O.

*2) The above applies to the field \(\mathbf{R}\) of real numbers. Every \(\mathbf{R}\) -algebra of finite degree that is a field is isomorphic to \(\mathbf{R}\) , \(\mathbf{C}\) , or \(\mathbf{H}\) .

\begin{enumerate}
\def\labelenumi{\arabic{enumi})}
\setcounter{enumi}{2}
\tightlist
\item
  Let A be a normed algebra over the field R. Suppose that A is a field. Then A is isomorphic to R, C, or H (``Gelfand--Mazur theorem'') (cf.~Comm. Alg., VI, §6, No.~4, p.~407, Theorem 1 and TS, I, §2, n° 5, p.~26, corollaire 2).*
\end{enumerate}

\BourbakiExercises{Exercises}

\begin{enumerate}
\def\labelenumi{\arabic{enumi})}
\tightlist
\item
  \begin{enumerate}
  \def\labelenumii{\alph{enumii})}
  \tightlist
  \item
    Let D be a noncommutative field with center K whose elements all have degree \(\leqslant 2\) over K. Prove that D is a quaternion algebra over K (Exercise 6, b) of VIII, p.~270).
  \end{enumerate}
\end{enumerate}

\begin{enumerate}
\def\labelenumi{\alph{enumi})}
\setcounter{enumi}{1}
\tightlist
\item
  Prove that a Cayley K-algebra that is a noncommutative field with center K is a quaternion algebra over K.
\end{enumerate}

¶ 2) a) Let A be a simple ring of finite rank over its center K, and let s be an involutive antiautomorphism of A. Let \(s'\) be an involutive antiautomorphism of A that coincides with s on K. Prove that there exists an element a of A such that we have \(s'(x) = as(x)a^{-1}\) for every \(x \in A\) and \(s(a) = a\) or \(s(a) = -a\) (apply the Skolem--Noether theorem and Theorem 3 of V, §11, No.~6, p.~85). Also prove the converse.

\begin{enumerate}
\def\labelenumi{\alph{enumi})}
\setcounter{enumi}{1}
\item
  Suppose that the restriction of s to K is not the identity; let \(K_{0}\) be the subfield of K consisting of the elements fixed by s. Prove that the set \(A_{0}\) (resp. \(A_{1}\) ) of elements a of A such that \(s(a) = a\) (resp. \(s(a) = -a\) ) is a \(K_{0}\) -linear subspace of A of dimension {[}A : K{]} and that \(A_{0}\) is a \(K_{0}\) -structure on A (use V, §10, No.~4, p.~63, Proposition 7).
\item
  Suppose that the restriction of s to K is the identity; set \([A:K]=m^{2}\) . Prove that \(A_{0}\) is a \(K_{0}\) -linear subspace of A of dimension \(m(m+1)/2\) or \(m(m-1)/2\) (reduce to the case when A is a matrix algebra over K and use a)).
\end{enumerate}

¶ 3) Let \((F, s)\) be a semisimple (associative) Cayley K-algebra.

\begin{enumerate}
\def\labelenumi{\alph{enumi})}
\item
  Prove that the K-algebra F is simple or isomorphic to the product of a simple algebra and its opposite algebra (consider the action of conjugation on the set of central idempotents).
\item
  Prove that if F is not commutative, then it is isomorphic to a quaternion algebra over K (use Exercise 1).
\item
  Prove that if F is commutative, then we are in one of the following cases:
\end{enumerate}

\(\alpha)\) F = K, and s is the identity.

\(\beta)\) F = K × K, and \(s(\lambda, \mu) = (\mu, \lambda)\) .

\(\gamma)\) F is a separable quadratic extension of K, and \(s\) is the nontrivial automorphism of F.

\(\delta)\) K has characteristic 2, F is an 2-radical extension of K of height 1, and \(s\) is the identity.

¶ 4) Let \((F, s)\) be an Artinian (associative) Cayley K-algebra, and let \(\mathfrak{r}\) be its radical. a) Prove that we have \(x^{2} = 0\) and \(s(x) = -x\) for every \(x \in \mathfrak{r}\) ; the algebra \(\overline{\mathrm{F}} = \mathrm{F} / \mathfrak{r}\) is Cayley and isomorphic to one of the algebras described in Exercise 3.

\begin{enumerate}
\def\labelenumi{\alph{enumi})}
\setcounter{enumi}{1}
\item
  Prove that if \(\overline{F}\) is a quaternion algebra over K, then \(\mathfrak{r}\) is reduced to 0 (observe that we have \((ab - ba)x = 0\) for \(a, b \in F, x \in \mathfrak{r}\) ).
\item
  Suppose that \(\overline{F}\) is commutative. Prove that \(\mathfrak{r}^{3}\) is reduced to 0 but that \(\mathfrak{r}^{2}\) may not be.
\item
  Suppose that \(\overline{F}\) is commutative and \(\mathfrak{r}^{2}\) is reduced to 0. Prove that there exists an \(\overline{F}\) -module V such that F is isomorphic to \(\overline{F} \oplus V\) endowed with the structure of a K-algebra for which we have \((\lambda, v)(\mu, w) = (\lambda \mu, \lambda w + s(\mu) v)\) . (In cases \(\alpha\) ) through \(\gamma\) ) of Exercise 3, c), use Corollary 1 of VIII, p.~245. In case \(\delta\) ), choose a p-basis \((e_{i})_{i \in I}\) of \(\overline{F}\) over K and for each \(i \in I\) , an element \(f_{i}\) of F that lifts \(e_{i}\) , and prove that the subalgebra of F generated by the \((f_{i})\) is isomorphic to \(\overline{F}\) .
\end{enumerate}

¶ 5) Suppose that the characteristic of K is ≠ 2. Let F be a not necessarily associative alternative Cayley K-algebra (III, Appendix, No.~2); we denote its conjugation by \(s : x \mapsto \overline{x}\) and its Cayley trace and norm by T and N, respectively. We say that a subalgebra E of F is nondegenerate if for every nonzero element x of E, there exists an \(x' \in E\) such that \(\mathrm{T}(x\overline{x}') \neq 0\) .

\begin{enumerate}
\def\labelenumi{\alph{enumi})}
\item
  Prove the equality \(x(\overline{z}y) + z(\overline{x}y) = \mathrm{T}(\overline{x}z)y = (yx)\overline{z} + (yz)\overline{x}\) for x, y, z in F (reduce to the case z = x).
\item
  Let E be a subalgebra of F containing 1, distinct from F, stable under s, and finite-dimensional over K; suppose that E and F are nondegenerate. Prove that there exists a nonzero element j of F such that \(T(xj) = 0\) for every \(x \in E\) and \(N(j) \neq 0\) . We have \(\overline{j} = -j\) , \(xj = -j\overline{x}\) for every \(x \in E\) , and \(j^{2} = \gamma1_{F}\) with \(\gamma = -N(j)\) .
\item
  Let \(x\) and \(y\) be elements of E. Prove the formulas \(x(yj) = (yx)j = (yj)\overline{x}\) and \((xj)(yj) = \gamma \overline{y}x\) (use \(a\) ) and Exercise 2 of III, Appendix, p.~654).
\end{enumerate}

\(d\) ) Prove that \(\mathrm{E} + \mathrm{E}j\) is a nondegenerate subalgebra of F, stable under \(s\) , and of dimension \(2\dim (\mathrm{E})\) that can be identified with the Cayley extension of \((\mathrm{E},s)\) defined by \(\gamma\) . Deduce from Proposition 3 of III, Appendix, No.~2, p.~614 that E is associative.

\begin{enumerate}
\def\labelenumi{\alph{enumi})}
\setcounter{enumi}{4}
\tightlist
\item
  Deduce from the above that the nondegenerate alternative Cayley K-algebras are the following:
\end{enumerate}

\begin{enumerate}
\def\labelenumi{(\roman{enumi})}
\item
  the K-algebra \(\mathrm{K}\) endowed with the identity involution
\item
  the quadratic algebras over K of type \((\alpha,0)\) with \(\alpha\neq0\)
\item
  the quaternion algebras over K of type \((\alpha,0,\gamma)\) with \(\alpha\gamma\neq0\)
\item
  the octonion algebras over \(\mathrm{K}\) of type \((\alpha, 0, \gamma, \delta)\) with \(\alpha \gamma \delta \neq 0\) .
\end{enumerate}

(Construct a sequence of Cayley extensions \(K \subset E_{1} \subset E_{2} \subset \cdots \subset F\) , and observe that an octonion algebra is not associative.)

\begin{enumerate}
\def\labelenumi{\alph{enumi})}
\setcounter{enumi}{5}
\tightlist
\item
  Prove that an alternative Cayley K-algebra whose nonzero elements are all invertible is one of the aforementioned types.
\end{enumerate}

\(g\) ) Let R a maximal ordered field. Prove that an alternative Cayley R-algebra whose nonzero elements are all invertible is isomorphic to R, C, H, or O (VIII, p.~368, Remark 1).

\begin{enumerate}
\def\labelenumi{\arabic{enumi})}
\setcounter{enumi}{5}
\tightlist
\item
  Let H be the quaternion field of type \((-1,0,-1)\) over Q.
\end{enumerate}

\begin{enumerate}
\def\labelenumi{\alph{enumi})}
\item
  An extension \(\mathrm{K}\) of \(\mathbf{Q}\) splits \(\mathrm{H}\) if and only if \(-1\) is the sum of two squares in \(\mathrm{K}\) (apply Propositions 3 of VIII, p.~363 and 4 of VIII, p.~364).
\item
  Let K be a cyclic extension of Q of degree \(2^{n}\) that splits H. Prove that no proper subfield of K splits H (observe that the unique subfield of K of degree \(2^{n-1}\) is the intersection of K and a maximal ordered field containing Q).
\item
  Let k be an integer \(\geqslant1\) and p a prime divisor of \(2^{2^{k}}+1\) . Prove that the greatest integer n such that \(2^{n}\) divides p-1 is \textgreater k (observe that we have \(2^{p-1}\equiv1\) modulo p). Let \(\Omega\) be a primitive p-the root of 1 and \(\mathrm{E}=\mathbf{Q}(\Omega)\) . Prove that -1 is the sum of two squares in E (let F be the extension \(\mathrm{E}(\sqrt{-1})\) ; show that -1 belongs to \(\mathrm{N}_{\mathrm{F}/\mathrm{E}}(\mathrm{F}^{*})\) by proving that this is true for \(\omega\) and \(1+\omega^{h}\) and by using the identity \(\prod_{j=0}^{r-1}(1+\mathrm{T}^{2^{j}})=\sum_{h=0}^{2^{r}-1}\mathrm{T}^{h})\) .
\end{enumerate}

\(d\) ) Prove that a cyclic extension of \(\mathbf{Q}\) of degree \(2^{n}\) contained in E splits H (apply VIII, p.~283, Corollary 2).

¶ 7) Suppose that the characteristic of K is different from 2.

\begin{enumerate}
\def\labelenumi{\alph{enumi})}
\item
  Let D be a field containing K and F the quaternion algebra over K of type \((\alpha, 0, \gamma)\) . Prove that \(D \otimes_{K} F\) is not a field if and only if there exist two commuting elements \(x\) and \(y\) in D such that \(x^2 - \alpha y^2 = \gamma\) (apply Exercise 4, c) of VIII, p.~351, first to \(D \otimes_{K} K(i)\) and then to \(D \otimes_{K} F\) ).
\item
  Let \(F'\) be another quaternion algebra over K. Then \(F \otimes_{K} F'\) is a field if and only if the only solution of \(\mathrm{N}_{\mathrm{F}}(q) = \mathrm{N}_{\mathrm{F}'}(q')\) , where \(q \in F\) and \(q' \in F'\) are pure quaternions (III, §2, p.~618, Exercise 3), is \(q = q' = 0\) . (Reduce to the case when F is a field and use the criterion of a). Note that if neither x nor y belongs to K, then y belongs to \(\mathrm{K}(x)\) , and write \(x = s + q\) , where \(x \in K\) and q is a pure quaternion.)
\end{enumerate}

\begin{enumerate}
\def\labelenumi{\arabic{enumi})}
\setcounter{enumi}{7}
\tightlist
\item
  Let \(K_{0}\) be a commutative field of characteristic different from 2, and let \((\mathrm{X}_{n})_{n\geqslant1}\) and \((\mathrm{Y}_{n})_{n\geqslant1}\) be two infinite sequences of variables and K be the field of rational fractions \(\mathrm{K}_{0}((\mathrm{X}_{n}),(\mathrm{Y}_{n}))\) . For any \(n\geqslant1\) , denote by \(F_{n}\) the field of quaternions over K of type \((\mathrm{X}_{n},0,\mathrm{Y}_{n})\) .
\end{enumerate}

\begin{enumerate}
\def\labelenumi{\alph{enumi})}
\item
  Prove that the tensor product \(F = \bigotimes_{n} F_{n}\) (III, §4, No.~5, p.~470) is a field with center K such that every subfield generated by a finite subset has finite rank over K (use Exercise 7, a)).
\item
  Prove that if u is an inner automorphism of F, then there exists a finite subset J of N such that the restriction of u to \(F_{n}\) , for \(n \notin J\) , is the identity. Deduce that there exists an uncountable set of noninner automorphisms of F.
\end{enumerate}

\section{Linear Representations of Algebras}

In this section, K is a commutative ring, and A is a K-algebra. From No.~2 on, we assume that K is a field, and we denote by \(A^{*}\) the dual of the vector space A over K. In this section, the notation \(A^{*}\) never refers to the multiplicative group of the ring A.

\subsection{Linear Representations of Algebras}

DEFINITION 1. --- Let M be a K-module. A linear representation of the algebra A in M is a K-homomorphism \(\pi\) from A to the algebra \(\operatorname{End}_{\mathrm{K}}(\mathrm{M})\) .

We also say that the pair \((\mathrm{M},\pi)\) is a linear representation of the algebra A. When M is a free K-module, the dimension of M is called the degree (or dimension) of \(\pi\) .

Let \(\pi\) be a linear representation of A in M. The additive law on M and the external law \((a,x)\mapsto\pi(a)(x)\) define a left A-module structure on M, which we denote by \(M_{\pi}\) . We say that \(M_{\pi}\) is the module of the representation \(\pi\) . The K-module structure on M is obtained from the A-module structure on \(M_{\pi}\) by restriction of scalars.

Conversely, let E be a left A-module. Let M be the K-module deduced from E by restriction of scalars from A to K, and let \(\pi\) be the homomorphism \(a \mapsto a_{E}\) from A to \(\operatorname{End}_{\mathrm{K}}(\mathrm{M})\) . Then \(\pi\) is a representation of A in M, and we have \(E = M_{\pi}\) . We say that \(\pi\) is the linear representation associated with the A-module E.

Studying A-modules or linear representations of A amounts to the same. We will translate certain definitions concerning modules to the language of linear representations.

Let \(\pi\) be a linear representation of A in M. The kernel of the homomorphism \(\pi\) is a two-sided ideal of A, which is simply the annihilator of the A-module \(\mathrm{M}_{\pi}\) . The homomorphism \(\pi\) is injective if and only if the A-module \(\mathrm{M}_{\pi}\) is faithful; we then say that \(\pi\) is a faithful representation of A.

Let \((\mathrm{M},\pi)\) and \((\mathrm{M}^{\prime},\pi^{\prime})\) be linear representations of A. An A-linear mapping from \(\mathrm{M}_{\pi}\) to \(\mathrm{M}_{\pi^{\prime}}^{\prime}\) is a K-linear mapping \(u:\mathrm{M}\to\mathrm{M}^{\prime}\) that satisfies the relation

\[
\pi^ {\prime} (a) \circ u = u \circ \pi (a)\tag{1}
\]

for \(a \in \mathbf{A}\) . An isomorphism from \(\mathrm{M}_{\pi}\) to \(\mathrm{M}_{\pi'}'\) is an isomorphism of K-modules \(u: \mathrm{M} \to \mathrm{M}'\) satisfying the condition

\[
\pi^ {\prime} (a) = u \circ \pi (a) \circ u ^ {- 1}\tag{2}
\]

for \(a \in A\) . We say that \(\pi\) and \(\pi'\) are isomorphic if the A-modules \(M_{\pi}\) and \(M'_{\pi'}\) are isomorphic. We say that \(\pi\) is a subrepresentation (resp. a quotient representation) of \(\pi'\) if \(M_{\pi}\) is a submodule (resp. a quotient module) of \(M'_{\pi'}\) .

Suppose given a family \((\mathrm{M}_{i},\pi_{i})\) of linear representations of A. Let M be the K-module that is the direct sum of the \(M_{i}\) ; for any \(a\in A\) , let \(\pi(a)\) be the endomorphism \((x_{i})_{i\in\mathrm{I}}\mapsto(\pi_{i}(a)(x_{i}))_{i\in\mathrm{I}}\) of M. Then \(\pi\) is a linear representation of A in M. The A-module \(M_{\pi}\) is the direct sum of the A-modules \((\mathrm{M}_{i})_{\pi_{i}}\) ( \(i\in I\) ); we say that \(\pi\) is the direct sum of the \(\pi_{i}\) , and we write \(\pi=\bigoplus\pi_{i}\) .

We say that the representation \(\pi\) of A in M is simple or irreducible if the A-module \(M_{\pi}\) is simple; we say that the representation \(\pi\) of A on M is semisimple or completely reducible if the A-module \(M_{\pi}\) is semisimple.

Let \(\pi\) be a linear representation of A in M. Let \(\mathrm{M}^*\) be the K-module dual to M. The mapping \(a\mapsto{}^{t}(\pi (a))\) from \(\mathrm{A}^{\circ}\) to \(\mathrm{End}_{\mathrm{K}}(\mathrm{M}^{*})\) is the transposed representation of \(\pi\) .

Let L be a commutative K-algebra, and let \((M,\pi)\) be a linear representation of the K-algebra A. The homomorphism \(\pi_{(L)}:A_{(L)}\to\operatorname{End}_{L}(M_{(L)})\) corresponding to the \(A_{(L)}\) -module structure on \(M_{(L)}\) is a linear representation of the L-algebra \(A_{(L)}\) . We say that \(\pi_{(L)}\) is the linear representation of the algebra \(A_{(L)}\) deduced from the representation \(\pi\) by extending the ring of scalars K to L.

Suppose that K is a field and that L is a nonzero commutative K-algebra.

Let \(\pi\) and \(\pi'\) be linear representations of the algebra A. It follows from VIII, p.~37, Theorem 3 that the representations \(\pi\) and \(\pi'\) are isomorphic if and only if \(\pi_{(L)}\) and \(\pi_{(L)}'\) are.

Suppose that K is a field. Let L be an extension of K. Consider the Grothendieck group \(R_{\mathrm{K}}(\mathrm{A})\) (resp. \(R_{\mathrm{L}}(\mathrm{A}_{(\mathrm{L})})\) ) of the A-modules that are finite-dimensional over K (resp. of the \(A_{(L)}\) -modules that are finite-dimensional over L). We have seen that the group homomorphism

\[
u: \mathrm{R} _ {\mathrm{K}} (\mathrm{A}) \longrightarrow \mathrm{R} _ {\mathrm{L}} (\mathrm{A} _ {(\mathrm{L})})
\]

defined by extension of scalars is injective; moreover, an element \(\xi \in \mathrm{R}_{\mathrm{K}}(\mathrm{A})\) is effective if and only if \(u(\xi)\) is (VIII, p.~195, Theorem 1).

\subsection{Restricted Dual of an Algebra}

We assume from now on that K is a field.

For \(a \in A\) , denote by \(\boldsymbol{\gamma}(a)\) the mapping \(x \mapsto ax\) and by \(\boldsymbol{\delta}(a)\) the mapping \(x \mapsto xa\) from A to A. Then \(\boldsymbol{\gamma}\) is a linear representation of A in A, called the left regular representation; likewise, \(\delta\) is a linear representation of \(A^{\circ}\) in A, called (by abuse of language) the right regular representation of A. By transposition, we deduce from \(\delta\) and \(\gamma\) linear representations of the algebras A and \(A^{\circ}\) in the vector space \(A^{*}\) that define on \(A^{*}\) a left A-module structure and a left \(A^{\circ}\) -module structure. These two structures correspond to an (A, A)-bimodule structure on \(A^{*}\) with external laws defined by the formulas

\[
\langle a f, b \rangle = \langle f, b a \rangle ,\tag{3}
\]

\[
\langle f a, b \rangle = \langle f, a b \rangle\tag{4}
\]

for \(a, b, \in \mathrm{A}\) and \(f \in \mathrm{A}^*\) .

Recall (II, §2, No.~4, p.~234, Definition 4) that if E is a linear subspace of A, then its orthogonal \(E'\) in \(A^{*}\) is the set of linear forms on A whose restriction to E is zero. Likewise, the orthogonal \(F'\) in A of a linear subspace F of \(A^{*}\) is the intersection of the kernels of the linear forms belonging to F.

We know (II, §7, No.~5, p.~299, Theorem 7) that the mapping \(\varphi : \mathrm{E} \mapsto \mathrm{E}'\) is a bijection from the set of linear subspaces of A of finite codimension to the set of linear subspaces of \(\mathrm{A}^*\) of finite dimension. The inverse mapping \(\varphi^{-1}\) sends a finite-dimensional subspace F of \(\mathrm{A}^*\) to its orthogonal \(\mathrm{F}'\) in A.

PROPOSITION 1. --- a) The mapping \(\varphi\) induces a bijection from the set of left (resp. right, two-sided) ideals of A of finite codimension to the set of right A-submodules (resp. left A-submodules, sub-bimodules) of \(A^{*}\) of finite dimension over K.

\begin{enumerate}
\def\labelenumi{\alph{enumi})}
\setcounter{enumi}{1}
\tightlist
\item
  Every left or right ideal of A of finite codimension contains a two-sided ideal of finite codimension.
\end{enumerate}

Formulas (3) and (4) prove that if E is a left ideal of A, then \(E'\) is a right A-submodule of \(A^{*}\) ; if V is a right A-submodule of \(A^{*}\) , then \(V'\) is a left ideal of A. The cases of right ideals and of two-sided ideals are treated analogously. This proves a).

Let \(\mathfrak{a}\) be a left ideal of A of finite codimension. The left A-module \(\mathrm{E} = \mathrm{A}_s / \mathfrak{a}\) is finite-dimensional over K, and so is \(\mathrm{End}_{\mathrm{K}}(\mathrm{E})\) . The kernel of the homomorphism \(a \mapsto a_{\mathrm{E}}\) from A to \(\mathrm{End}_{\mathrm{K}}(\mathrm{E})\) is therefore a two-sided ideal of A of finite codimension contained in \(\mathfrak{a}\) . Passing to the opposite ring \(\mathrm{A}^{\circ}\) , we see that every right ideal of A of finite codimension contains a two-sided ideal of finite codimension.

DEFINITION 2. --- The restricted dual

of the K-algebra A, denoted by \(\Theta(A)\) , is the union in \(A^{*}\) of the orthogonals of the two-sided ideals of \(A\) of finite codimension.

By Proposition 1, we can give the following equivalent descriptions of \(\Theta(A)\) :

- the union of the orthogonals of the left ideals of A of finite codimension

-- the union of the orthogonals of the right ideals of A of finite codimension

- the union of the orthogonals of the two-sided ideals of \(\mathrm{A}\) of finite codimension

- the union of the left A-submodules of \(\mathrm{A}^*\) of finite dimension over \(\mathbf{K}\)

- the union of the right A-submodules of \(\mathrm{A}^*\) of finite dimension over \(K\)

- the union of the \((\mathrm{A},\mathrm{A})\) -sub-bimodules of \(\mathrm{A}^*\) of finite dimension over \(\mathbf{K}\) .

We have \(\Theta(A) = \Theta(A^{\circ})\) , and \(\Theta(A) = A^{*}\) if \(A\) has finite degree over \(K\) . Let \(f \in A^{*}\) ; then \(f\) belongs to \(\Theta(A)\) if and only if \(Af\) (resp. \(fA\) , \(AfA\) ) is a linear subspace of \(A^{*}\) of finite dimension over \(K\) .

The sum of two (A, A)-sub-bimodules of \(A^{*}\) of finite dimension over K has finite dimension over K. It follows that \(\Theta(\mathrm{A})\) is an (A, A)-sub-bimodule of \(A^{*}\) .

\subsection{Coefficients of a Module}

Let E be a left A-module. Let \(E^{*}\) be the dual of the vector space E over K; we endow it with its natural right A-module structure. Given elements x of E and \(x^{*}\) of \(E^{*}\) , we denote by \(c_{\mathrm{E}}(x, x^{*})\) the linear form

\[
c _ {\mathrm{E}} (x, x ^ {*}): a \mapsto \langle x ^ {*}, a x \rangle\tag{5}
\]

on A. These linear forms are called the coefficients of E. The subspace of \(\mathrm{A}^*\) generated by the coefficients of E is denoted by \(\Theta_{\mathrm{E}}(\mathrm{A})\) . It is nonzero if E is nonzero.

If F is an A-module isomorphic to E, then E and F have the same coefficients and we have \(\Theta_{\mathrm{F}}(\mathrm{A}) = \Theta_{\mathrm{E}}(\mathrm{A})\) . Let us view \(E^{*}\) as a left \(A^{\circ}\) -module; a coefficient of E is a coefficient of \(E^{*}\) , and \(\Theta_{\mathrm{E}}(\mathrm{A}) \subset \Theta_{\mathrm{E}^{*}}(\mathrm{A}^{\circ})\) . Consequently, if E is finite-dimensional over K, then E and \(E^{*}\) have the same coefficients, and we have \(\Theta_{\mathrm{E}}(\mathrm{A}) = \Theta_{\mathrm{E}^{*}}(\mathrm{A}^{\circ})\) .

Let \((\mathrm{M},\pi)\) be a linear representation of the algebra A. The coefficients of the A-module \(\mathrm{M}_{\pi}\) are also called the coefficients of the representation \(\pi\) . The coefficient \(c_{\mathrm{M}_{\pi}}(x,x^{*})\) , for \(x\in \mathrm{M}\) and \(x^{*}\in \mathrm{M}^{*}\) , is also denoted by \(c_{\pi}(x,x^{*})\) ; the vector space \(\Theta_{\mathrm{M}_{\pi}}(\mathrm{A})\) is also denoted by \(\Theta_{\pi}(\mathrm{A})\) .

Remarks. --- 1) Suppose that M is finite-dimensional over K. Let \((e_1, \ldots, e_n)\) be a basis of V, and let \((e_1^*, \ldots, e_n^*)\) be its dual basis. Denote by \((\pi_{ij}(a))\) the matrix of \(\pi(a)\) with respect to the basis \((e_1, \ldots, e_n)\) of V; we have \(\pi_{ij} = c_\pi(e_j, e_i^*)\) . The mapping \(a \mapsto (\pi_{ij}(a))\) is a homomorphism from A to \(\mathbf{M}_n(\mathrm{K})\) ; such a homomorphism is sometimes called a matrix representation of the algebra A.

\begin{enumerate}
\def\labelenumi{\arabic{enumi})}
\setcounter{enumi}{1}
\tightlist
\item
  Let E be an A-module, and let \(E'\) be a submodule of E. By the corollary of Theorem 5 of II, §7, No.~5, p.~299, we have \(\Theta_{\mathrm{E}'}(\mathrm{A}) \subset \Theta_{\mathrm{E}}(\mathrm{A})\) and \(\Theta_{\mathrm{E}/\mathrm{E}'}(\mathrm{A}) \subset \Theta_{\mathrm{E}}(\mathrm{A})\) .
\end{enumerate}

Let \(\gamma_{\mathrm{E}}\) be the unique K-linear mapping from \(\mathrm{E} \otimes_{\mathrm{K}} \mathrm{E}^{*}\) to \(\mathrm{A}^{*}\) that sends \(x \otimes x^{*}\) to \(c_{\mathrm{E}}(x, x^{*})\) . It is (A,A)-bilinear, and its image is \(\Theta_{\mathrm{E}}(\mathrm{A})\) . We deduce from it A-linear mappings \(c_{\mathrm{E}}^{\prime}: \mathrm{E} \to \operatorname{Hom}_{\mathrm{A}}(\mathrm{E}^{*}, \mathrm{A}^{*})\) and \(c_{\mathrm{E}}^{\prime \prime}: \mathrm{E}^{*} \to \operatorname{Hom}_{\mathrm{A}}(\mathrm{E}, \mathrm{A}^{*})\) such that

\[
c _ {\mathrm{E}} ^ {\prime} (x) (x ^ {*}) = c _ {\mathrm{E}} (x, x ^ {*}) = c _ {\mathrm{E}} ^ {\prime \prime} (x ^ {*}) (x).
\]

Denote by \(\theta_{E}\) the K-linear mapping \(\theta_{E}: E \otimes E^{*} \to \operatorname{End}_{K}(E)\) characterized by the relation

\[
\theta_ {\mathrm{E}} (x \otimes x ^ {*}) (y) = \langle x ^ {*}, y \rangle x
\]

for \(x, y \in E\) and \(x^{*} \in E^{*}\) . Its image is the set \(\operatorname{End}_{\mathrm{K}}^{\mathrm{f}}(\mathrm{E})\) of endomorphisms of E of finite rank (VIII, p.~463). By the definition of the trace (loc. cit.), we have

\[
\operatorname{Tr} \left(\theta_ {\mathrm{E}} \left(x \otimes x ^ {*}\right)\right) = \langle x ^ {*}, x \rangle
\]

for \(x \in E\) and \(x^{*} \in E^{*}\) ; we therefore have

\[
\langle \gamma_ {\mathrm{E}} (x \otimes x ^ {*}), a \rangle = \langle x ^ {*}, a x \rangle = \operatorname{Tr} (\theta_ {\mathrm{E}} (a x \otimes x ^ {*})) = \operatorname{Tr} (\theta_ {\mathrm{E}} (x \otimes x ^ {*}) a).
\]

This gives the relation

\[
\langle \gamma_ {\mathrm{E}} (h), a \rangle = \operatorname{Tr} (\theta_ {\mathrm{E}} (h) \circ a _ {\mathrm{E}})
\]

for \(a \in \mathrm{A}\) and \(h \in \mathrm{E} \otimes_{\mathrm{K}} \mathrm{E}^*\) . This proves that \(\Theta_{\mathrm{E}}(\mathrm{A})\) is the set of linear forms \(a \mapsto \operatorname{Tr}(u \circ a_{\mathrm{E}})\) on \(\mathrm{A}\) , where \(u\) runs through \(\operatorname{End}_{\mathrm{K}}^{\mathrm{f}}(\mathrm{E})\) .

Lemma 1. --- The mapping \(c_{\mathrm{E}}^{\prime \prime}\) is bijective. If \(\mathrm{E}\) is finite-dimensional, then so is \(c_{\mathrm{E}}^{\prime}\) .

When F is a right A-module and G is a K-vector space, we defined in II, §4, No.~1, p.~268, Proposition 1, a) an isomorphism \(\gamma\) of K-vector spaces from \(\mathrm{Hom}_{\mathrm{K}}(\mathrm{F} \otimes_{\mathrm{A}} \mathrm{E}, \mathrm{G})\) to \(\mathrm{Hom}_{\mathrm{A}}(\mathrm{E}, \mathrm{Hom}_{\mathrm{K}}(\mathrm{F}, \mathrm{G}))\) . This isomorphism sends \(\varphi : \mathrm{F} \otimes_{\mathrm{A}} \mathrm{E} \to \mathrm{G}\) to the homomorphism \(e \mapsto (f \mapsto \varphi(f \otimes e))\) . Through the canonical isomorphism from \(\mathrm{A}_d \otimes_{\mathrm{A}} \mathrm{E}\) to E (II, §3, No.~4, p.~249), \(\gamma\) can be identified with \(c_{\mathrm{E}}''\) when \(\mathrm{F} = \mathrm{A}_d\) and \(\mathrm{G} = \mathrm{K}\) .

Analogously, the isomorphism \(\beta\) defined in II, §4, No.~1, p.~268, Proposition 1, b) specializes to an isomorphism \(\beta\) from \(\mathrm{Hom}_{\mathrm{K}}(\mathrm{E}^{*}\otimes_{\mathrm{A}}\mathrm{A}_{s},\mathrm{K})\) to \(\mathrm{Hom}_{\mathrm{A}}(\mathrm{E}^{*},\mathrm{Hom}_{\mathrm{K}}(\mathrm{A}_{s},\mathrm{K}))\) . When E is finite-dimensional, E can be canonically identified with \(\mathrm{Hom}_{\mathrm{K}}(\mathrm{E}^{*},\mathrm{K})\) and \(\mathrm{E}^{*}\) with \(\mathrm{E}^{*}\otimes_{\mathrm{A}}\mathrm{A}_{s}\) ; the homomorphism \(\beta\) is then identified with \(c_{\mathrm{E}}^{\prime}\) .

PROPOSITION 2. --- Let E be a left A-module.

\begin{enumerate}
\def\labelenumi{\alph{enumi})}
\item
  The set of coefficients of E is the union of the images of the A-linear mappings from E to A*.
\item
  Suppose, moreover, that E has finite dimension n over K. Then the left A-module \(\Theta_{\mathrm{E}}(\mathrm{A})\) is isomorphic to a quotient of \(E^{n}\) , the A-module E is isomorphic to a submodule of \(\Theta_{\mathrm{E}}(\mathrm{A})^{n}\) , and every element of \(\Theta_{\mathrm{E}}(\mathrm{A})\) is a coefficient of \(E^{n}\) .
\end{enumerate}

Assertion a) follows from the surjectivity of \(c_{\mathrm{E}}^{\prime \prime}\) .

Let us prove b). Let \((e_1^*,\ldots ,e_n^*)\) be a basis of \(\mathrm{E}^*\) over K. Since \(\Theta_{\mathrm{E}}(\mathrm{A})\) is generated by the coefficients of E, the A-linear mapping

\[
(x _ {1}, \ldots , x _ {n}) \mapsto \sum_ {i = 1} ^ {n} c _ {\mathrm{E}} (x _ {i}, e _ {i} ^ {*})
\]

from \(\mathrm{E}^n\) to \(\Theta_{\mathrm{E}}(\mathrm{A})\) is surjective. By a), every element of \(\Theta_{\mathrm{E}}(\mathrm{A})\) is a coefficient of \(\mathrm{E}^n\) . Moreover, the A-linear mapping

\[
x \mapsto \left(c _ {\mathrm{E}} (x, e _ {1} ^ {*}), \ldots , c _ {\mathrm{E}} (x, e _ {n} ^ {*})\right)
\]

from E to \(\Theta_{\mathrm{E}}(\mathrm{A})^{n}\) is injective; this gives b).

Remark 3. --- Let E be a left A-submodule of \(\mathrm{A}^*\) . Let \(\varepsilon\) be the linear form \(y \mapsto y(1)\) on E. For every \(x\) in E, we have \(x = c_{\mathrm{E}}(x, \varepsilon)\) , so E is an A-submodule of \(\Theta_{\mathrm{E}}(\mathrm{A})\) .

\subsection{Restricted Dual and Matrix Coefficients}

PROPOSITION 3. --- Let \((\mathrm{V},\pi)\) be a finite-dimensional linear representation of the algebra A. The kernel \(\mathfrak{a}\) of \(\pi\) is a two-sided ideal of A of finite codimension, and \(\Theta_{\pi}(\mathrm{A})\) is the orthogonal of \(\mathfrak{a}\) in \(A^{*}\) . The transpose of the mapping \(\pi\) defines an isomorphism from the dual of the K-vector space \(\pi(\mathrm{A})\) to \(\Theta_{\pi}(\mathrm{A})\) .

Since \(\operatorname{End}_{\mathrm{K}}(\mathrm{V})\) is finite-dimensional over K, \(\mathfrak{a}\) is a two-sided ideal of A of finite codimension. The vector space \(\Theta_{\pi}(\mathrm{A})\) is a finite-dimensional subspace of \(A^{*}\) , and its orthogonal in A is equal to \(\mathfrak{a}\); by Theorem 7 of II, §7, No.~5, p.~301, the space \(\Theta_{\pi}(\mathrm{A})\) is therefore the orthogonal of \(\mathfrak{a}\) in \(A^{*}\) . Moreover, since \(\mathfrak{a}\) is the kernel of \(\pi\) , the transpose of the mapping \(\pi\) defines an isomorphism from the dual of \(\pi(\mathrm{A})\) to the orthogonal of \(\mathfrak{a}\) in \(A^{*}\) (II, §7, No.~5, p.~299, Corollary of Theorem 5).

COROLLARY. --- The restricted dual \(\Theta(\mathrm{A})\) of A is the set of coefficients of the finite-dimensional linear representations of A.

By definition, \(\Theta(A)\) is the union of the orthogonals of the two-sided ideals of \(A\) of finite codimension. We therefore have \(\Theta_{\pi}(A) \subset \Theta(A)\) for every finite-dimensional linear representation \(\pi\) of \(A\) . Conversely, let \(f\) be an element of \(\Theta(A)\) , and let \(\mathfrak{a}\) be a two-sided ideal of finite codimension, contained in the kernel of \(f\) (VIII, p.~376, Definition 2). Denote by \(\pi\) the linear representation of \(A\) in \(A/\mathfrak{a}\) deduced from the left regular representation in \(A\) by passing to the quotient. Let \(x\) be the class of 1 (mod \(\mathfrak{a}\) ), and let \(x^{*}\) be the linear form on \(A/\mathfrak{a}\) deduced from \(f\) . We have \(f = c_{\pi}(x, x^{*})\) , so that \(f\) is a coefficient of \(\pi\) .

\pandocbounded{\includegraphics[keepaspectratio]{images/part003_35b408b4.jpg}}

Take note that if \((V,\pi)\) is a linear representation of A that is not finite-dimensional over K, then the space \(\Theta_{\pi}(A)\) is not necessarily contained in \(\Theta(A)\) .

\subsection{Dual of a Semisimple Algebra}

Let \(\Theta^{\mathrm{ss}}(\mathrm{A})\) be the socle of the left A-module \(\Theta (\mathrm{A})\) , that is, (VIII, p.~65) the largest semisimple submodule of \(\Theta (\mathrm{A})\) . We denote by \(\mathcal{S}_{\mathrm{K}}\) the set of classes of simple (left) A-modules that are finite-dimensional over K. When A is a semisimple algebra of finite degree over K, we have \(\mathrm{A}^{*} = \Theta (\mathrm{A}) = \Theta^{\mathrm{ss}}(\mathrm{A})\) because every left A-module is semisimple (VIII, p.~138, Proposition 4).

THEOREM 1. --- a) The set \(\Theta^{\mathrm{ss}}(\mathrm{A})\) consists of the coefficients of the finite-dimensional semisimple representations of A. It is an (A,A)-sub-bimodule of \(\mathrm{A}^*\) .

\begin{enumerate}
\def\labelenumi{\alph{enumi})}
\setcounter{enumi}{1}
\item
  For every \(S \in \mathcal{S}_{\mathrm{K}}\) , the isotypical component of \(\Theta(A)\) of type S is equal to \(\Theta_{\mathrm{S}}(\mathrm{A})\) . The left A-module \(\Theta^{\mathrm{ss}}(\mathrm{A})\) is the direct sum of the submodules \(\Theta_{\mathrm{S}}(\mathrm{A})\) , where S runs through \(\mathcal{S}_{\mathrm{K}}\) .
\item
  For every S in \(\mathcal{S}_{\mathrm{K}}\) , the right A-module \(S^{*}\) is simple, and \(\Theta_{\mathrm{S}}(\mathrm{A})\) is the isotypical component of type \(S^{*}\) of the right A-module \(\Theta(\mathrm{A})\) . The mapping that sends S to \(\operatorname{cl}(\mathrm{S}^{*})\) is a bijection from \(\mathcal{S}_{\mathrm{K}}\) to the set of classes of simple \(A^{o}\) -modules of finite dimension over K.
\item
  Viewed as a right A-module, \(\Theta (\mathrm{A})\) has socle \(\Theta^{\mathrm{ss}}(\mathrm{A})\) .
\end{enumerate}

Let E be a semisimple A-module of finite dimension over K. Every coefficient of E belongs to the image of an A-linear mapping from E to \(A^{*}\) (VIII, p.~378, Proposition 2) and therefore belongs to \(\Theta^{\mathrm{ss}}(\mathrm{A})\) . Conversely, let \(f \in \Theta^{\mathrm{ss}}(\mathrm{A})\) . Then the A-module Af is finite-dimensional over K and is semisimple. By the remark in VIII, p.~367, f is a coefficient of Af.

For every S in \(\mathcal{S}_{\mathrm{K}}\) , the isotypical component of \(\Theta(\mathrm{A})\) of type S is generated by the images of the A-linear mappings from S to \(A^{*}\) ; it is therefore equal to \(\Theta_{\mathrm{S}}(\mathrm{A})\) by Proposition 2, a) of VIII, p.~378. On the other hand, if S is a simple A-module that is not finite-dimensional over K, then the isotypical component of \(\Theta(\mathrm{A})\) of type S is zero: indeed, every simple submodule of \(\Theta(\mathrm{A})\) is monogenous and therefore finite-dimensional over K. Since the socle of \(\Theta(\mathrm{A})\) is the direct sum of its isotypical components (VIII, p.~65, Proposition 4), this proves b).

Let S be a simple A-module of finite dimension over K. Since the K-vector space S is not reduced to 0, the same holds for \(S^{*}\) . Let E be a submodule of the right A-module \(S^{*}\) ; its orthogonal \(E'\) in S is an A-submodule of S. Since S is simple, we have either \(E' = 0\) , in which case \(E = S^{*}\) , or \(E' = S\) , in which case E = 0. Hence \(S^{*}\) is a simple right A-module.

We have \(\Theta (\mathrm{A}^{\circ}) = \Theta (\mathrm{A})\) (VIII, p.~376); we identify right A-modules with left \(\mathrm{A}^{\circ}\) -modules. Since every vector space of finite dimension over K is isomorphic to its bidual, the above proves that the mapping \(\mathrm{S}\mapsto \mathrm{cl}(\mathrm{S}^{*})\) is a bijection from \(\mathcal{S}_{\mathrm{K}}\) to the set of classes of simple \(\mathrm{A}^{\circ}\) -modules of finite dimension over K. Now, for S in \(\mathcal{S}_{\mathrm{K}}\) , the isotypical component of \(\Theta (\mathrm{A}^{\circ})\) of type \(\mathrm{S}^*\) is equal to \(\Theta_{\mathrm{S}^*}(\mathrm{A}^{\circ})\) by assertion b) applied to the algebra \(\mathrm{A}^{\circ}\) , and we have \(\Theta_{\mathrm{S}^*}(\mathrm{A}^{\circ}) = \Theta_{\mathrm{S}}(\mathrm{A})\) . Assertions c) and d) follow immediately.

COROLLARY 1. --- Every simple A-module of finite dimension over K is isomorphic to an A-submodule of \(\Theta(A)\) .

This follows from Theorem 1, b) because we have \(\Theta_{\mathrm{E}}(\mathrm{A}) \neq 0\) for every nonzero A-module E.

COROLLARY 2. --- If two simple A-modules of finite dimension over K have a common nonzero coefficient, then they are isomorphic.

Let \(S_{1}\) and \(S_{2}\) be simple A-modules of finite dimension over K with a common nonzero coefficient. We then have \(\Theta_{S_{1}}(A) \cap \Theta_{S_{2}}(A) \neq 0\) . Now, \(\Theta_{S_{1}}(A)\) is isotypical of type \(S_{1}\) , and \(\Theta_{S_{2}}(A)\) is isotypical of type \(S_{2}\) . So \(S_{1}\) is isomorphic to \(S_{2}\) .

By Theorem 1, \(\Theta^{\mathrm{ss}}(\mathrm{A})\) is an (A,A)-sub-bimodule of \(\mathrm{A}^*\) , the direct sum of the (A,A)-bimodules \(\Theta_{\mathrm{S}}(\mathrm{A})\) for S running through \(\mathcal{S}_{\mathrm{K}}\) . These bimodules are pairwise nonisomorphic because they are already nonisomorphic as left A-modules.

Fix an S in \(\mathcal{S}_{\mathrm{K}}\) . Denote by D the opposite field of \(\operatorname{End}_{\mathrm{A}}(\mathrm{S})\) , and view S as a right D-module and \(S^{*}\) as a left D-module. With these conventions, S is an (A, D)-bimodule, \(S^{*}\) is a (D, A)-bimodule, and \(S \otimes_{D} S^{*}\) is an (A, A)-bimodule.

PROPOSITION 4. --- a) There exists a group homomorphism \(\lambda_{\mathrm{S}}\) from \(\mathrm{S} \otimes_{\mathrm{D}} \mathrm{S}^{*}\) to \(\Theta_{\mathrm{S}}(\mathrm{A})\) , characterized by

\[
\lambda_ {\mathrm{S}} (x \otimes x ^ {*}) = c _ {\mathrm{S}} (x, x ^ {*})\tag{6}
\]

for \(x \in S\) and \(x^{*} \in S^{*}\) . This mapping is an isomorphism of (A, A)-bimodules.

\begin{enumerate}
\def\labelenumi{\alph{enumi})}
\setcounter{enumi}{1}
\tightlist
\item
  The (A, A)-bimodule \(\Theta_{\mathrm{S}}(\mathrm{A})\) is simple.
\end{enumerate}

The mapping \(\gamma_{\mathrm{S}}: \mathrm{S} \otimes_{\mathrm{K}} \mathrm{S}^{*} \to \Theta_{\mathrm{S}}(\mathrm{A})\) characterized by the formula \(\gamma_{\mathrm{S}}(x \otimes x^{*}) = c_{\mathrm{S}}(x, x^{*})\) is (A, A)-linear and surjective and satisfies \(\gamma_{\mathrm{S}}(xd \otimes x^{*}) = \gamma_{\mathrm{S}}(x \otimes dx^{*})\) for \(x \in \mathrm{S}\) , \(x^{*} \in \mathrm{S}^{*}\) , and \(d \in \mathrm{D}\) . It therefore defines a surjective (A, A)-linear mapping \(\lambda_{\mathrm{S}}: \mathrm{S} \otimes_{\mathrm{D}} \mathrm{S}^{*} \to \Theta_{\mathrm{S}}(\mathrm{A})\) characterized by formula (6).

Let us prove that \(S \otimes_{D} S^{*}\) is a simple (A, A)-bimodule. By Corollary 2 of VIII, p.~63, every (A, A)-sub-bimodule of \(S \otimes_{D} S^{*}\) is of the form \(S \otimes_{D} H\) , where H is a (D, A)-sub-bimodule of \(S^{*}\) . Since \(S^{*}\) is a simple right A-module (VIII, p.~380, Theorem 1, c)), we have H = 0 or \(H = S^{*}\) ; the assertion follows.

The homomorphism \(\lambda_{S}\) is (A, A)-linear and nonzero; since \(S \otimes_{D} S^{*}\) is a simple bimodule, \(\lambda_{S}\) is injective (VIII, p.~47, Proposition 2, a)).

Remark. --- When the field K is algebraically closed, we have D = K by Theorem 1 of VIII, p.~47, and \(\lambda_{S}\) is an isomorphism of (A, A)-bimodules from \(S \otimes_{K} S^{*}\) to \(\Theta_{\mathrm{S}}(\mathrm{A})\) .
